\documentclass[10pt,a4paper]{amsart}
\usepackage[margin=19mm]{geometry}
\usepackage{amsmath,amssymb,mathtools}
\usepackage[scr=rsfs,cal=euler]{mathalfa}
\usepackage{xfrac}
\usepackage[unicode,hidelinks,bookmarks=false]{hyperref}
\newcommand{\ie}{i.e.\ }
\newcommand{\cf}{cf.\ }
\newcommand{\resp}{respectively\ }
\newcommand{\Subsection}{\subsection}
\providecommand{\nobreakdash}{\nobreak\hbox{-}}
\setlength{\emergencystretch}{2em}
\multlinegap0pt
\newcommand{\skpt}{\par\smallskip}
\providecommand{\eg}{e.g.\ }
\usepackage{enumerate}
\usepackage{esint}
\usepackage{aligned-overset}
\usepackage{upgreek}
\usepackage{cleveref}
\crefdefaultlabelformat{\emph{#1}}
\crefname{step}{\emph{Step}}{Steps}

\newcommand{\R}{\mathbb R}
\newcommand{\C}{\mathbb C}
\newcommand{\bbZ}{\mathbb Z}

\DeclareMathOperator{\Id}{Id}
\DeclareMathOperator{\esssup}{esssup}
\DeclareMathOperator{\Tr}{Tr}
\DeclareMathOperator{\Cof}{Cof}
\DeclareMathOperator{\Cov}{Cov}
\DeclareMathOperator{\Div}{div}
\DeclareMathOperator{\Vol}{Vol}
\DeclareMathOperator{\locc}{loc}

\newcommand{\Pheat}{\mathrm{P}}
\newcommand{\Hheat}{\mathrm{H}}
\newcommand{\map}{\mathrm T}
\newcommand{\Ent}{\mathrm H}
\newcommand{\KM}{\mathrm F}


\newcommand{\V}{V}
\newcommand{\W}{W}
\newcommand{\sta}{g}
\newcommand{\fin}{h}
\newcommand{\dd}{n}
\newcommand{\conv}{c}
\newcommand{\density}{f}

\newcommand{\prob}{\rho}
\newcommand{\source}{\mu}
\newcommand{\target}{\nu}
\newcommand{\lsh}{\alpha}
\newcommand{\lc}{\kappa}

\newcommand*\diff{\mathop{}\!\mathrm{d}}

\newcommand{\cotrans}{\mathcal L}
\newcommand{\clN}{\mathcal{N}}
\newcommand{\clZ}{\mathcal{Z}}

\newcommand{\Gaussian}{\upgamma}
%\newcommand{\Gaussianw}{\psi}
\newcommand{\OT}{\Phi}
\newcommand{\wave}{\psi}
\newcommand{\waveg}{\psi_{\mathrm{Gaussian}}}
\newcommand{\cxfun}{\varphi}
\newcommand{\avgint}{\fint}
\newcommand{\findiff}{\delta\OT}
%%%%%%%%%%%%%%%%%%%%%%%%%%%%%%%%%%%%%%%%%%%%%%%%
\renewcommand*{\to}{\mathchoice{\longrightarrow}{\rightarrow}{\rightarrow}{\rightarrow}}
%%%
\let\oldmapsto\mapsto
\renewcommand*{\mapsto}{\mathchoice{\longmapsto}{\oldmapsto}{\oldmapsto}{\oldmapsto}}
%%%%
\renewcommand*{\implies}{\mathchoice{\Longrightarrow}{\Rightarrow}{\Rightarrow}{\Rightarrow}}


\let\originalleft\left 
\let\originalright\right
\renewcommand{\left}{\mathopen{}\mathclose\bgroup\originalleft}
\renewcommand{\right}{\aftergroup\egroup\originalright}

\let\oldin\in
\renewcommand{\in}{\mathchoice{\oldin}{\oldin}{\,\oldin\,}{\,\oldin\,}}

\newcounter{setofstep}
\newcounter{step}[setofstep]
\newenvironment{step}[1][]{\refstepcounter{step}\begin{proof}[{Step~\thestep\if#1\empty\else: #1\fi}]}{\let\qed\relax\end{proof}}
\newcommand{\newstepset}{\stepcounter{setofstep}}

\newcommand*{\mk}{\mkern -1mu}
\newcommand*{\Mk}{\mkern -2mu}
\newcommand*{\mK}{\mkern 1mu}
\newcommand*{\MK}{\mkern 2mu}


\newcommand*{\relabel}{\renewcommand{\labelenumi}{(\theenumi)}}
\newcommand*{\romanenumi}{\renewcommand*{\theenumi}{\roman{enumi}}\relabel}
\newcommand*{\Romanenumi}{\renewcommand*{\theenumi}{\Roman{enumi}}\relabel}
\newcommand*{\alphenumi}{\renewcommand*{\theenumi}{\alph{enumi}}\relabel}
\newcommand*{\Alphenumi}{\renewcommand*{\theenumi}{\Alph{enumi}}\relabel}
\let\oldtilde\tilde
\renewcommand*{\tilde}[1]{\mathchoice{\widetilde{#1}}{\widetilde{#1}}{\oldtilde{#1}}{\oldtilde{#1}}}
\let\oldforall\forall
\renewcommand*{\forall}{\mathrel{\oldforall}}

\newtheorem{theo}{Theorem}[section]
\newtheorem{lemm}[theo]{Lemma}
\newtheorem{prop}[theo]{Proposition}
\theoremstyle{definition}\newtheorem{defi}[theo]{Definition}
\newtheorem{rema}[theo]{Remark}
\newtheorem{assumption}[theo]{Assumption}
\newenvironment{enonce}[1]{\begin{assumption}}{\end{assumption}}
\newenvironment{proofc}[1][]{\begin{proof}[#1]}{\end{proof}}
\title[Brenier Laplacian bound]{Optimal transport maps, majorization, and log-subharmonic measures: original Theorem2.1}
\author{Guido De Philippis and Yair Shenfeld}
\date{Source: Annales Henri Lebesgue8(2025),925--963; DOI10.5802/ahl.251}
\begin{document}\maketitle
\noindent\textit{Editorial scope.} The counted unit is original Theorem2.1 with its complete proof and required smooth-approximation and AppendixA local core. The full-support and curvature assumptions remain exact. Later applications are omitted. Original print slips, including the literal $\lim V=0$ after the preceding decay estimate, the empty $L^{}$ exponent and finite-difference notation, are recorded without generated corrections; every original construction and bound remains.
\section{Original definitions and Lipschitz observation}
\begin{defi}\label{def:log_concave}
A probability measure $\prob$ on $\R^{\dd}$, with density $\diff \prob=e^{-U}\diff x$ with respect to the Lebesgue measure, is \emph{$\conv$-log-convex} (res. \emph{$\conv$-log-concave}), for $\conv\in \R$, if the distributional derivative $\nabla^2U$ satisfies $\nabla^2U \preceq \conv \Id_{\dd}$ (res. $\nabla^2U \succeq \conv \Id_{\dd}$), where $\preceq$ (res. $\succeq$) are in the sense of positive semidefinite order.
\end{defi}


\begin{theo}[Caffarelli~\cite{MR1800860}, Kolesnikov~\cite{kolesnikov2011mass}]\label{thm:Caffarelli_intro}
Let $\diff\source=e^{-\V}\diff x$ and $\diff\target=e^{-\W}\diff x$ be probability measures on $\R^{\dd}$, with $\source$ supported on all of $\R^{\dd}$, such that there exist $\lsh>0,\lc>0$ with
\[
\nabla^2 \V \preceq \lsh \Id_{\dd}\quad\text{and}\quad \nabla^2\W \succeq \lc \Id_{\dd}.
\]
Let $\nabla\OT:\R^{\dd}\to\R^{\dd}$ be the Brenier map transporting $\source$ to $\target$. Then,
\begin{equation}\label{eq:Caffarelli_intro}
\left\|\nabla^2\OT\right\|_{L^{\infty}(\diff x)}\le \sqrt{ \frac{\lsh}{\lc}},
\end{equation}
where $\|\cdot\|_{L^{\infty}(\diff x)}$ is the $L^{\infty}$ operator norm.
\end{theo}

\setcounter{theo}{2}
\begin{defi}\label{def:log_subharmonic}
A probability measure $\prob$ on $\R^{\dd}$, with density $\diff \prob=e^{-U}\diff x$ with respect to the Lebesgue measure where $U\in L_{\locc}^1(\diff x)$, is \emph{$\conv$-log-subharmonic} (res. \emph{$\conv$-log-superharmonic}), for $\conv\in \R$, if the distributional derivative $\Delta U$ satisfies $\Delta U \le\conv $ (res. $\Delta U \ge \conv$).
\end{defi}


\setcounter{theo}{4}
\begin{rema}[Lipschitz bounds]\label{rem:Lipschiz}
Since the Brenier map $\nabla\OT$ between $\source$ and $\target$ is a gradient of a convex function, our bound~\href{https://ahl.centre-mersenne.org/articles/10.5802/ahl.251/}{(1.4)} implies that under the assumptions of Theorem~\href{https://ahl.centre-mersenne.org/articles/10.5802/ahl.251/}{1.4} we get the Lipschitz bound
\begin{equation}\label{eq:Lipschitz}
\left\|\nabla^2\OT\right\|_{L^{\infty}(\diff x)}\le\dd \sqrt{ \frac{\lsh}{\lc}}.
\end{equation}
The bound~\eqref{eq:Lipschitz} is sharp as can be seen by taking
\[
\source=\clN(0,\Sigma_{\epsilon}),\quad \target=\clN(0,\Id_{\dd}); \quad \lc=1, \quad \lsh_{\epsilon}=1+\frac{\dd-1}{\dd}\epsilon,
\]
where $\Sigma_{\epsilon}$ is a diagonal matrix with one entry equal to $1/\dd$ and the rest of the entries equal to $1/\epsilon$ for some $\epsilon\in (0,1]$. In this case $\nabla\OT(x)=\Sigma_{\epsilon}^{-1}x$ so that
\[
\left\|\nabla^2\OT\right\|_{L^{\infty}(\diff x)}\le \dd=\lim_{\epsilon\to 0}\dd \sqrt{ \frac{\lsh_{\epsilon}}{\lc}}.
\]
The bound~\eqref{eq:Lipschitz} can be seen as the dimensional price that one has to pay if the log-convexity assumption in Theorem~\href{https://ahl.centre-mersenne.org/articles/10.5802/ahl.251/}{1.2} is relaxed to log-subharmonicity assumption.
\end{rema}

\section{The divergence of the Brenier map}\label{sec:divergence_OT}
In this section we will prove $L^p$ estimates on $\Delta\OT$ (Theorem~\ref{thm:main_OT_Lp}), which in particular will imply Theorem~\href{https://ahl.centre-mersenne.org/articles/10.5802/ahl.251/}{1.4}. While the proof of our main result below Theorem~\ref{thm:main_OT_Lp} is quite technical, at the formal level it follows the original (formal) derivation of Caffarelli~\cite{MR1800860}, as we now explain. Our goal is to derive the bound
\begin{equation}\label{eq:infty_bound_main_sec}
\|\Delta\OT\|_{L^{\infty}(\diff x)}\le \dd\sqrt{ \frac{\lsh}{\lc}} \quad \text{when}\quad \Delta \V\le \lsh\dd\quad\text{and}\quad \nabla^2\W\succeq \lc \Id_{\dd}.
\end{equation}
The derivation is based on differentiating the Monge--Ampère equation twice, and using the optimality conditions at the point where $\Delta\OT$ attains its maximum. To this end, given a unit vector $e\in \R^{\dd}$ and a function $\xi:\R^{\dd}\to \R$, denote by $\xi_e$ (res. $\xi_{ee}$) the first (res. second) directional derivative of $\xi$ in the direction $e$. The Monge--Ampère equation reads
\begin{equation}\label{eq:Monge}
e^{-\V}=e^{-\W}(\nabla \OT)\det\nabla^2\OT,
\end{equation}
so take the logarithm on both sides of~\eqref{eq:Monge}, and differentiate twice in a fixed direction $e$, to get
\begin{align}\label{eq:Vee}
\V_{ee}&=\left\langle \nabla^2\W(\nabla\OT)\nabla^2\OT e,\nabla^2\OT e\right\rangle+\langle \nabla\W,\nabla\OT_{ee}\rangle-\partial_{ee}\log\det \nabla^2\OT.
\end{align}
Using the identity
\begin{equation}\label{eq:Jacobi_ee}
\partial_{ee}\log\det \nabla^2\OT=\Tr\left[(\nabla^2\OT)^{-1}\nabla^2\OT_{ee}\right]-\Tr\left[\left((\nabla^2\OT)^{-1}\nabla^2\OT_e\right)^2\right],
\end{equation}
the identity~\eqref{eq:Vee} becomes
\begin{multline}\label{eq:Vee+}
\V_{ee}=\left\langle \nabla^2\W(\nabla\OT)\nabla^2\OT e,\nabla^2\OT e\right\rangle+\langle \nabla\W,\nabla\OT_{ee}\rangle\\
-\Tr\left[(\nabla^2\OT)^{-1}\nabla^2\OT_{ee}\right]+\Tr\left[\left((\nabla^2\OT)^{-1}\nabla^2\OT_e\right)^2\right].
\end{multline}
Using $\nabla^2\W\succeq \lc \Id_{\dd}$, \eqref{eq:Vee+} implies
\begin{multline}\label{eq:Vee_inq}
\V_{ee}\ge \lc\left|\nabla^2\OT e\right|^2+\langle \nabla\W,\nabla\OT_{ee}\rangle\\
-\Tr\left[(\nabla^2\OT)^{-1}\nabla^2\OT_{ee}\right]
+\Tr\left[\left((\nabla^2\OT)^{-1}\nabla^2\OT_e\right)^2\right],
\end{multline}
so summing on both sides of~\eqref{eq:Vee_inq} over a basis $\{e_i\}$ of $\R^{\dd}$ yields
\begin{multline}\label{eq:Jacobi_Laplacian}
\Delta \V\ge\lc \sum_{i=1}^{\dd}\left|\nabla^2\OT e_i\right|^2+ \langle \nabla\W,\nabla\Delta\OT\rangle\\
-\Tr\left[(\nabla^2\OT)^{-1}\nabla^2\Delta\OT\right]+\sum_{i=1}^{\dd}\Tr\left[\left((\nabla^2\OT)^{-1}\nabla^2\OT_{e_i}\right)^2\right].
\end{multline}



Suppose $\Delta\OT$ attains its maximum $x_0$. Then the optimality conditions give $\nabla \Delta\OT(x_0)=0$ and $\nabla^2\Delta\OT(x_0) \preceq 0$, so
\begin{equation}\label{eq:opt_condition}
\langle \nabla\W(x_0),\nabla\Delta\OT(x_0)\rangle-\Tr\left[\left(\nabla^2\OT(x_0)\right)^{-1}\nabla^2\Delta\OT(x_0)\right]\ge 0.
\end{equation}
On the other hand,
\begin{equation}\label{eq:Tr_noneg}
\Tr\left[\left((\nabla^2\OT)^{-1}\nabla^2\OT_e\right)^2\right]=\Tr[A^2]\ge 0\quad\text{with}\quad A:=(\nabla^2\OT)^{-\frac{1}{2}}\nabla^2\OT_e(\nabla^2\OT)^{-\frac{1}{2}}.
\end{equation}
Hence, combining~\eqref{eq:Jacobi_Laplacian}, \eqref{eq:opt_condition}, \eqref{eq:Tr_noneg}, and applying the Cauchy--Schwarz inequality, shows that
\begin{equation}\label{eq:main_inq_proof_e}
\lsh\dd\ge\Delta \V(x_0)\ge \lc \sum_{i=1}^{\dd}\left|\nabla^2\OT(x_0) e_i\right|^2\ge \kappa\sum_{i=1}^{\dd}\left\langle \nabla^2\OT(x_0)e_i,e_i\right\rangle^2.
\end{equation}
By Jensen's inequality,
\begin{equation}\label{eq:main_inq_proof}
\lsh\dd\ge\Delta \V(x_0)\ge \lc\sum_{i=1}^{\dd}\left\langle \nabla^2\OT(x_0)e_i,e_i\right\rangle^2\ge \frac{\lc}{\dd}[\Delta\OT(x_0)]^2,
\end{equation}
so it follows that
\begin{equation}\label{eq:OT_subharmonic}
\Delta\OT(x_0)\le \dd\sqrt{\frac{\lsh}{\lc}},
\end{equation}
where we used that $\Delta\OT(x_0)\ge 0$ since $\OT$ is convex. Since $x_0$ is a point where $\Delta\OT$ attains its maximum, we conclude~\eqref{eq:infty_bound_main_sec}. Making the above argument rigorous is difficult due to the need to show the existence of a point $x_0$ where $\Delta\OT$ attains its maximum. Instead, we follow the $L^p$ approach of Kolesnikov~\cite{MR3201654}, \cite[\S~6]{kolesnikov2011mass}, but with some necessary modifications (cf. Remark~\ref{rem:modification}).




\begin{theo}\label{thm:main_OT_Lp}
Let $\diff\source=e^{-\V}\diff x$ and $\diff\target=e^{-\W}\diff x$ be probability measures on~$\R^{\dd}$, with $\source$ supported on all of $\R^{\dd}$, such that there exist $\lsh>0,\lc>0$ with
\[
\Delta \V \le \lsh\dd \quad\text{and}\quad \nabla^2\W \succeq \lc \Id_{\dd}.
\]
Let $\nabla\OT:\R^{\dd}\to\R^{\dd}$ be the Brenier map transporting $\source$ to $\target$. Then,
\begin{align}\label{eq:infty_bound_main}
\|\Delta\OT\|_{L^{\infty}(\diff x)}\le \dd\sqrt{ \frac{\lsh}{\lc}}.
\end{align}
\end{theo}

We start by showing that it suffices to prove Theorem~\ref{thm:main_OT_Lp} assuming sufficient regularity for $\source$ and $\target$. These regularity assumptions are captured as follows.

\begin{enonce}{Assumption}\label{assump}\leavevmode
\begin{enumerate}
\item\label{Assumpt2.2.1} $\V$ is smooth.
\item\label{Assumpt2.2.2} There exists a constant $c>0$ such that $\nabla^2\V(x) \preceq c\Id_{\dd}$ for all $x\in \R^{\dd}$.
\item\label{Assumpt2.2.3} There exist constants $a,b >0$ such that $\lim_{|x|\to\infty}\V(x)-(b|x|^2-a)\ge 0$.
\item\label{Assumpt2.2.4} $\W$ is smooth.
\item\label{Assumpt2.2.5} There exists a constant $d>0$ such that $\nabla^2\W(x) \preceq d\Id_{\dd}$ for all $x\in \R^{\dd}$.
\end{enumerate}
\end{enonce}

Note that Assumption~\ref{assump}$\MK$\eqref{Assumpt2.2.3} implies that $\source$ has a finite second moment while the assumption $\nabla^2\W \succeq \lc \Id_{\dd}$ implies that $\target$ has a finite second moment (by the Poincaré inequality for $\target$).

\begin{prop}\label{prop:smooth}
It suffices to prove Theorem~\ref{thm:main_OT_Lp} for $\source,\target$ which in addition to satisfying the assumptions of Theorem~\ref{thm:main_OT_Lp} also satisfy Assumption~\ref{assump}.
\end{prop}

\begin{proof}
The proof is based on the composition of the following two steps:
\begin{itemize}%[left=0pt]
\item[Step~\ref{prfprop2.4step1}:] Given $\source,\target$ which satisfy the assumptions of Theorem~\ref{thm:main_OT_Lp} we show that there exist sequences of probability measures $\{\source_k\},\{\target_k\}$ converging weakly to $\source,\target$, respectively, such that, for all $k$, each pair $\source_k,\target_k$ satisfies the assumptions of Theorem~\ref{thm:main_OT_Lp} and Assumption~\ref{assump}.

\item[Step~\ref{prfprop2.4step2}:] Assume that Theorem~\ref{thm:main_OT_Lp} holds true for each pair $\source_k,\target_k$ and then pass to the limit to deduce~\eqref{eq:infty_bound_main}.
\end{itemize}

Before we start with Step~\ref{prfprop2.4step1} let us state a number of important smoothing properties of the Ornstein--Uhlenbeck semigroup which we will use to smooth out the measures $\source$ and $\target$.


\begin{prop}\label{prop:OU_smooth}
Let $\Gaussian$ be the standard Gaussian measure in $\R^{\dd}$, and let $f:\R^{\dd}\to \R_{\ge 0}$ be a nonnegative function in $L^1(\Gaussian)$. Let $(\Pheat_t)_{t\ge 0}$ be the Ornstein--Uhlenbeck semigroup,
\begin{equation}\label{eq:OU_def_prop}
\Pheat_t\density(x):=\int_{\R^{\dd}} \density\left(e^{-t}x+\sqrt{1-e^{-2t}}y\right)\diff \Gaussian(y),\quad t\ge 0,\quad x\in \R^{\dd}.
\end{equation}
\begin{enumerate}\romanenumi
\item\label{prop2.4.1} For every $x\in \R^{\dd}$ and $t>0$,
\begin{equation}\label{eq:less_lc}
\nabla^2\log\Pheat_t\density(x) \succeq -\frac{e^{-2t}}{1-e^{-2t}} \Id_{\dd}\quad \textnormal{for all }x\in\R^{\dd} \text{ and }t>0.
\end{equation}
\item\label{prop2.4.2} Suppose $\density$ is $\conv$-log-concave for $\conv\in \R$ (i.e., $x\mapsto \log \density(x)-\conv\frac{|x|^2}{2}$ is concave). Then, for every $x\in \R^{\dd}$,
\begin{equation}\label{eq:lc_OU}
\nabla^2\log\Pheat_t\density(x)\preceq \frac{ e^{-2t}\conv}{1-\conv(1-e^{-2t})} \Id_{\dd}
\begin{cases}
\text{for any }t\in [0,\infty) &\text{if }\conv\le 1\\
\text{for any }t\in \left[0,\log\left(\sqrt{\frac{\conv}{\conv-1}}\right)\right] &\text{if }\conv >1.
\end{cases}
\end{equation}
\item\label{prop2.4.3} Suppose $\density$ is $\conv\dd$-log-subharmonic for some $\conv\le 0$ (i.e., $x\mapsto \log f(x)-\conv\frac{|x|^2}{2}$ is subharmonic). Then, for every $x\in \R^{\dd}$ and $t>0$,
\begin{equation}\label{eq:beta_subharmonic}
\Delta\log\Pheat_t\density(x)\ge \frac{e^{-2t}\conv\dd}{1-\conv\left(1-e^{-2t}\right)}\ge \conv\dd.
\end{equation}
\end{enumerate}
\end{prop}

The proof of Proposition~\ref{prop:OU_smooth} (with some extra results) is given in Section~\ref{subsec:appendix_smooth}.


Let us begin with Step~\ref{prfprop2.4step1}.

\begin{step}\label{prfprop2.4step1}
We first construct the sequence $\{\source_k\propto e^{-\V_k}\diff x\}$ whose members all satisfy the assumptions of Theorem~\ref{thm:main_OT_Lp} as well as Assumptions~\ref{assump}$\MK$\eqref{Assumpt2.2.1}--\eqref{Assumpt2.2.3}. Given $\source=e^{-\V}\diff x$ let
\[
\tilde V_k:=-\log\Pheat_{\frac{1}{k}}e^{-\V}.
\]
Then $\tilde V_k$ is smooth, and by Proposition~\ref{prop:OU_smooth}$\MK$\eqref{prop2.4.3} it satisfies $\Delta\tilde\V_k(x)\le \lsh\dd$ for all $x\in \R^{\dd}$. This shows that we can construct a measure $\propto e^{-\tilde \V_k}\diff x$ which satisfies the assumptions of Theorem~\ref{thm:main_OT_Lp} as well as Assumption~\ref{assump}$\MK$\eqref{Assumpt2.2.1}. Moreover, Assumption~\ref{assump}$\MK$\eqref{Assumpt2.2.2} is also satisfied for some $c_k$ by Proposition~\ref{prop:OU_smooth}$\MK$\eqref{prop2.4.1}. Next we modify $\tilde V_k$ so that Assumption~\ref{assump}$\MK$\eqref{Assumpt2.2.3} is also satisfied. Define
\[
\V_k(x):=\left(1-\frac{1}{k}\right)\tilde \V_k(x)+\frac{1}{k}\lsh\frac{|x|^2}{2}
\]
and note that $\Delta\V_k\le \lsh\dd$, $\V_k$ is smooth, and
\[
\nabla^2\V_k \preceq \left(\left(1-\frac{1}{k}\right)c_k+\frac{1}{k}\lsh\right)\Id_{\dd},
\]
so the measure $\propto e^{- \V_k}\diff x$ satisfies the assumptions of Theorem~\ref{thm:main_OT_Lp} as well as Assumptions~\ref{assump}$\MK$\eqref{Assumpt2.2.1}--\eqref{Assumpt2.2.2}. Let us show that $\V_k$ also satisfies Assumption~\ref{assump}$\MK$\eqref{Assumpt2.2.3}. To this end we first need to argue that $\V$ is positive outside of some ball. Indeed, since $\Delta\V\le \lsh\dd$, we have that, for any $x_0\in \R^{\dd}$, the function
\[
\R^{\dd}\ni x\mapsto e^{\lsh\frac{|x-x_0|^2}{2}-\V(x)}
\]
is subharmonic. Hence, choosing $x_0$ such that $|x_0|>R$ for some $R>0$, we have
\begin{equation}\label{eq:V_bd}
e^{-\V(x_0)}\le \avgint_{B_1(x_0)}e^{\lsh\frac{|x-x_0|^2}{2}-\V(x)}\diff x\le \frac{e^{\frac{\lsh}{2}}}{\Vol(B_1)}\int_{(B_{R-1}(0))^c}e^{-\V(x)}\diff x.
\end{equation}
The right-hand side of~\eqref{eq:V_bd} converges to $0$ as $R\to\infty$ since $\int e^{-V}$ is finite. Hence, $\lim_{x\to\infty}\V(x)=0$ and Assumption~\ref{assump}$\MK$\eqref{Assumpt2.2.3} follows by the construction of $\V_k$. Finally, let $\source_k \propto e^{-\V_k}\diff x$ and note that $\source_k\to \source$ weakly.

Next we construct the sequence $\{\target_k\propto e^{-\W_k}\diff x\}$ whose members all satisfy the assumptions of Theorem~\ref{thm:main_OT_Lp} as well as Assumptions~\ref{assump}$\MK$\eqref{Assumpt2.2.4}--\eqref{Assumpt2.2.5}. Given $\target=e^{-\W}\diff x$ let
\[
\W_k:=-\log\Pheat_{\frac{1}{k}}e^{-\W}
\]
and let $\{\target_k\propto e^{-\W_k}\diff x\}$. Then, as in the construction of $\V_k$, $\W_k$ is smooth and there exists a constant $d_k$ such that $\nabla^2\W \preceq d_k\Id_{\dd}$. Hence, $\W_k$ satisfies Assump\-tions~\ref{assump}\linebreak\eqref{Assumpt2.2.4}--\eqref{Assumpt2.2.5}. Finally, by Proposition~\ref{prop:OU_smooth}$\MK$\eqref{prop2.4.2}, $\W_k$ satisfies the assumptions of Theorem~\ref{thm:main_OT_Lp} with $\nabla^2\W_k \succeq \lc_k$ where
\begin{equation}\label{eq:lc_k}
\lc_k:= \frac{\lc e^{-\frac{2}{k}}}{1+\lc\left(1-e^{-\frac{2}{k}}\right)} =\lc-\frac{\lc(\lc +1)\left(1-e^{-\frac{2}{k}}\right)}{1+\lc\left(1-e^{-\frac{2}{k}}\right)}.
\end{equation}
\end{step}
We now move to Step~\ref{prfprop2.4step2}. 
\begin{step}\label{prfprop2.4step2}
Assume that Theorem~\ref{thm:main_OT_Lp} holds true for each pair $\source_k,\target_k$, and let $\nabla\OT_k$ be the Brenier map between $\source_k$ and $\target_k$. Then, for $k$ large enough, by Theorem~\ref{thm:main_OT_Lp} and Remark~\ref{rem:Lipschiz},
\begin{equation}\label{eq:Delta_Phi_k}
\|\Delta\OT_k\|_{L^{\infty}(\diff x)}\le \dd\sqrt{ \frac{\lsh}{\lc_k}}\quad\text{and}\quad \nabla^2\OT_k(x) \preceq \dd\sqrt{ \frac{\lsh}{\lc_k}}\Id_{\dd} \text{ for every $x\in \R^{\dd}$}.
\end{equation}
By~\cite[Lemma~1]{MS2022}, which follows~\cite[Lemma~2.1]{neeman2022lipschitz} building on~\cite[Lem\-ma~3.3]{MR2983070}, we get that, up to a subsequence, $\{\nabla\OT_k\}$ converges almost everywhere to some transport map $T$ between $\source$ and~$\target$. Since $\{\nabla\OT_k\}$ converges so does $\{\OT_k\}$ to some convex function $\OT$. It follows that $T=\nabla\OT$ is the Brenier map between $\source$ and~$\target$. Finally, the proof is complete since \(\Phi_{k}\) converges to \(\Phi\) locally uniformly, \(\Delta \Phi_{k}\) converges to \(\Delta \Phi\) as distribution, and as
\begin{align*}
\|\Delta \Phi\|_{L^{\infty}(\diff x)}&=\sup_{\substack{\eta \in C_{c}^{\infty},\|\eta\|_{L^1}=1}} \int \Delta \Phi(x) \eta(x) \diff x =\sup_{\substack{\eta \in C_{c}^{\infty}, \|\eta\|_{L^1}=1}} \int \Phi(x) \Delta \eta(x) \diff x
\\
&=\sup_{\substack{\eta \in C_{c}^{\infty}, \|\eta\|_{L^1}=1}} \lim_{k \to \infty} \int \Phi_{k}(x) \Delta \eta(x) \diff x\\
&=\sup_{\eta \in C_{c}^{\infty}, \|\eta\|_{L^{}(\diff x)}=1} \lim_{k \to \infty} \int \Delta \Phi_{k}(x) \eta(x) \diff x
\\
&\le \lim_{k\to \infty} \|\Delta \Phi_{k}\|_{L^{\infty}(\diff x)}=\dd\sqrt{ \frac{\lsh}{\lc}}.\qedhere
\end{align*}
\end{step}\let\qed\relax
\end{proof}

\begin{proof}[Proof of Theorem~\ref{thm:main_OT_Lp}]
In light of Proposition~\ref{prop:smooth} we may assume Assumption~\ref{assump} from here on. Under the assumptions of Theorem~\ref{thm:main_OT_Lp}, together with Assumption~\ref{assump}, we have that there exists an optimal transport map $\nabla\OT$ between $\source$ and $\target$ which is smooth~\cite[Theorem~1.1]{Dario2019}, satisfies the Monge--Ampère equation,
\begin{equation}\label{eq:Monge_proof}
e^{-\V(x)}=e^{-\W(\nabla \OT(x))}\det\nabla^2\OT(x), \quad \text{for all }x\in \R^{\dd},
\end{equation}
and, by Theorem~\ref{thm:Caffarelli_intro}, satisfies the bound
\begin{equation}\label{eq:Caffarelli_bound}
0\preceq\nabla^2\OT(x)\preceq C\Id_{\dd}\quad \text{for all }x\in \R^{\dd},
\end{equation}
for some constant $C>0$. Taking the logarithm in~\eqref{eq:Monge_proof} we get
\begin{align}\label{eq:Monge_log}
\V(x)=\W(\nabla\OT(x))-\log \det \nabla^2\OT(x),
\end{align}
so, for every $x,y\in \R^{\dd}$,
\begin{multline}\label{eq:Monge_finite_diff}
V(x+y)+V(x-y)-2V(x)\\
=\bigl\{\W(\nabla\OT(x+y))+\W(\nabla\OT(x-y))-2\W(\nabla\OT(x))\bigr\}
\\
-\log\left[\frac{\det \nabla^2\OT(x+y)\det \nabla^2\OT(x-y)}{(\det \nabla^2\OT(x))^2}\right].
\end{multline}
Equation~\eqref{eq:Monge_finite_diff} is the finite difference analogue of equation~\eqref{eq:Vee}. The left-hand side of~\eqref{eq:Monge_finite_diff} is a second-order finite difference of $\V$ which we wish to relate to $\Delta\V$, so to this end we introduce the following $\Delta_{\epsilon}$ operator and its properties; see Section~\ref{subsec:appendix_Delta_eps} for the proof.


\begin{lemm}\label{lem:avg_Laplacian}
Fix $\epsilon> 0$ and for a function $f:\R^{\dd}\to\R$ denote
\begin{equation}\label{eq:avg_Laplacian_def}
\Delta_{\epsilon}f(x):=\avgint_{\partial B_{\epsilon}(0)}[f(x+y)-f(x)]\diff y,
\end{equation}
where
\[
\avgint_{\partial B_{\epsilon}(0)}f:=\frac{1}{|\partial B_{\epsilon}(0)|}\int_{\partial B_{\epsilon}(0)}f.
\]
Then,
\begin{align}\label{eq:Delta_eps_lim}
\lim_{\epsilon\to 0}\frac{\Delta_{\epsilon}f(x)}{\epsilon^{2}}=\frac{\Delta f(x)}{2\dd},
\end{align}
where $\Delta f$ is the distributional Laplacian of $f$. Further, if $\Delta f\le \ell$, then
\begin{align}\label{eq:Delta_eps_bound}
\Delta_{\epsilon}f\le \frac{\ell}{\dd}\frac{\epsilon^2}{2},\quad \quad\forall~\epsilon>0.
\end{align}
\end{lemm}

We now integrate~\eqref{eq:Monge_finite_diff} over $y$ with respect to the uniform measure on the sphere $\partial B_{\epsilon}(x)$ of radius $\epsilon$ centered at $x$. By the assumption $\Delta V\le \lsh\dd$, together with~\eqref{eq:Delta_eps_bound}, we have
\begin{multline}\label{eq:Monge_avg}
\epsilon^2 \lsh\ge 2\Delta_{\epsilon}V(x)=2\avgint_{\partial B_{\epsilon}(0)}\bigl[\W(\nabla\OT(x+y))-\W(\nabla\OT(x))\bigr]\diff y
\\
-\avgint_{\partial B_{\epsilon}(0)}\log\left[\frac{\det \nabla^2\OT(x+y)\det \nabla^2\OT(x-y)}{(\det \nabla^2\OT(x))^2}\right]\diff y.
\end{multline}
On the other hand, by the assumption $\nabla^2\W \succeq \lc \Id_{\dd}$,
\begin{multline}\label{eq:cvx_W}
\W(\nabla\OT(x+y))-\W(\nabla\OT(x))\\
\ge \bigl\langle (\nabla \W)(\nabla\OT(x)),\nabla\OT(x+y)-\nabla\OT(x)\bigr\rangle+\frac{\lc}{2}|\nabla\OT(x+y)-\nabla\OT(x)|^2,
\end{multline}
which is the analogue of~\eqref{eq:Vee_inq}. Combining~\eqref{eq:Monge_avg} and~\eqref{eq:cvx_W} implies
\begin{multline}\label{eq:Monge_avg_cvx}
\epsilon^2 \lsh \ge \lc\avgint_{\partial B_{\epsilon}(0)}|\nabla\OT(x+y)-\nabla\OT(x)|^2\diff y
\\
+\avgint_{\partial B_{\epsilon}(0)}\bigg\{2\bigl\langle (\nabla \W)(\nabla\OT(x)),\nabla\OT(x+y)-\nabla\OT(x)\bigr\rangle\\
\left.-\log\left[\frac{\det \nabla^2\OT(x+y)\det \nabla^2\OT(x-y)}{(\det \nabla^2\OT(x))^2}\right]\right\}\diff y.
\end{multline}
Next we will multiply both sides of~\eqref{eq:Monge_avg_cvx} by $(\Delta_{\epsilon}\OT(x))^p$ and integrate against $e^{-V(x)}\diff x$. This requires showing that some integrals are finite.

\begin{lemm}\label{lem:integrate_finite}
For every $p>0$ we have
\begin{multline}\label{eq:2_terms}
\epsilon^2 \lsh\int (\Delta\OT_{\epsilon}(x))^p e^{-V(x)}\diff x\\
\shoveleft{\ge \lc\int\avgint_{\partial B_{\epsilon}(0)}\left[|\nabla\OT(x+y)-\nabla\OT(x)|^2 \right]\diff y(\Delta_{\epsilon}\OT(x))^p e^{-V(x)}\diff x}\\
%+\int (\Delta_{\epsilon}\OT(x))^p e^{-V(x)}\cdot\left\{\avgint_{\partial B_{\epsilon}(0)}\bigg\{ 2\bigl\langle \nabla\W(\nabla \OT(x)),\nabla \OT(x+y)-\nabla \OT(x)\bigl\rangle\right.\\
\left.\left.-\log\left[\frac{\det \nabla^2\OT(x+y)\det \nabla^2\OT(x-y)}{(\det \nabla^2\OT(x))^2}\right]\right\}\diff y\right\}\diff x,
\end{multline}
and all the integrals are finite.
\end{lemm}
\begin{proof}
Once we show that the integrals are finite, Equation~\eqref{eq:2_terms} follows by multiplying both sides of~\eqref{eq:Monge_avg_cvx} by $(\Delta_{\epsilon}\OT)^p$ and integrating against $e^{-V}\diff x$. Note that at the moment we are only interested in these quantities to be finite and not uniformly bounded in \(\epsilon\). The integral $\int (\Delta_{\epsilon}\OT)^p e^{-V}\diff x$ is finite by~\eqref{eq:Caffarelli_bound} and~\eqref{eq:Delta_eps_bound}. The integral $\int\avgint_{\partial B_{\epsilon}(0)}\left[|\nabla\OT(x+y)-\nabla\OT(x)|^2 \right]\diff y(\Delta_{\epsilon}\OT)^p e^{-V}\diff x$ is finite since by the fundamental theorem of calculus we can write the difference $\nabla\OT(x+y)-\nabla\OT(x)$ as an integral along a path between $x+y$ and $x$ of $\nabla^2\OT$ multiplied by a vector of length $\le\epsilon$, so the difference is bounded since $\nabla^2\OT$ is bounded by~\eqref{eq:Caffarelli_bound}. To bound the term $\langle \nabla\W(\nabla \OT(x)),\nabla \OT(x+y)-\nabla \OT(x)\rangle$ use Assumption~\ref{assump}$\MK$\eqref{Assumpt2.2.5} according to which $\nabla W$ grows at most linearly, and the fact that again $\nabla \OT(x+y)-\nabla \OT(x)$ can be bounded by a constant, to see that $\int (\Delta_{\epsilon}\OT)^p\avgint_{\partial B_{\epsilon}(0)} 2\langle \nabla\W(\nabla \OT(x)),\nabla \OT(x+y)-\nabla \OT(x)\rangle e^{-V}$ is finite since $\target$ has a finite first moment. Finally, to bound the term
\[
\int\avgint_{\partial B_{\epsilon}(0)}\left| \log\left[\frac{\det \nabla^2\OT(x+y)\det \nabla^2\OT(x-y)}{(\det \nabla^2\OT(x))^2}\right]\right| e^{-V(x)}\diff x \diff y
\]
we use identity~\eqref{eq:Monge_avg} so we need to bound
\[
\int\avgint_{\partial B_{\epsilon}(0)}\bigl|\W(\nabla\OT(x+y))-\W(\nabla\OT(x))\bigr|(\Delta_{\epsilon}\OT)^p e^{-V}\diff y\diff x
\]
as well as
\[
\int |\Delta_{\epsilon}\V| e^{-V}\diff x.
\]
For the first term, again by~\eqref{eq:Caffarelli_bound}, the distance between $\nabla\OT(x+y)$ and $\nabla\OT(x)$ is bounded by a constant, so by Assumption~\ref{thm:main_OT_Lp}$\MK$\eqref{Assumpt2.2.5}, $\avgint_{\partial B_{\epsilon}(0)}|\W(\nabla\OT(x+y))-\W(\nabla\OT(x))|$ is linear in $\nabla\OT$ (at some point between $x+y$ and $x$) and, since $\nabla^2\Phi$ is bounded, $\int\avgint_{\partial B_{\epsilon}(0)}|\W(\nabla\OT(x+y))-\W(\nabla\OT(x))|(\Delta_{\epsilon}\OT)^p e^{-V}\diff y\diff x$ is finite since $\target$ has a finite first moment. Finally, for the term $ \int |\Delta_{\epsilon}\V |e^{-V}\diff x$, we use the definition of $\Delta_{\epsilon}$ and note that at infinity $\V$ grows quadratically so $\int |\Delta_{\epsilon}\V |e^{-V}\diff x$ is finite since $\source$ has a finite second moment.
\end{proof}

We will show in Proposition~\ref{prop:noneg_term} below that the second term in~\eqref{eq:2_terms} is nonnegative, which is the replacement of~\eqref{eq:opt_condition}--\eqref{eq:Tr_noneg}. Let us complete the proof of Theorem~\ref{thm:main_OT_Lp} assuming the validity of Proposition~\ref{prop:noneg_term}.
\begin{lemm}\label{lem:_term_to_thm_proof}
Suppose that for every $p>0$,
\begin{multline}\label{eq:1_term}
\epsilon^2 \lsh\int (\Delta_{\epsilon}\OT(x))^p e^{-V(x)}\diff x\\
\ge \lc\int\avgint_{\partial B_{\epsilon}(0)}\left[|\nabla\OT(y+x)-\nabla\OT(x)|^2 \right]\diff y(\Delta_{\epsilon}\OT(x))^p e^{-V(x)}\diff x.
\end{multline}
Then,
\begin{align}\label{eq:infty_bound}
\|\Delta\OT\|_{L^{\infty}(\source)}\le \dd\sqrt{ \frac{\lsh}{\lc}}.
\end{align}
\end{lemm}
\begin{proof}
Dividing both sides of~\eqref{eq:1_term} by $\epsilon^{2p+2}$, and using~\eqref{eq:Delta_eps_lim}, we get
\begin{multline*}
\lsh \int \left(\frac{\Delta \OT}{2\dd}\right)^p e^{-V}\diff x\\
\begin{aligned}
&\ge \lc \int\left\{\lim_{\epsilon\to 0}\frac{1}{\epsilon^2}\avgint_{\partial B_{\epsilon}(0)}\left[|\nabla\OT(y+x)-\nabla\OT(x)|^2 \right]\diff y\right\} \left(\frac{\Delta \OT}{2\dd}\right)^p e^{-V}\diff x\\
&=\lc \int\left\{\avgint_{\partial B_{1}(0)}\lim_{\epsilon\to 0}\left|\frac{\nabla\OT(x+\epsilon y)-\nabla\OT(x)}{\epsilon}\right|^2\diff y\right\} \left(\frac{\Delta \OT)}{2\dd}\right)^p e^{-V}\diff x\\
&=\lc \int\left\{\avgint_{\partial B_{1}(0)}\left| \nabla^2\OT(x)y\right|^2\diff y\right\} \left(\frac{\Delta \OT}{2\dd}\right)^p e^{-V}\diff x.
\end{aligned}
\end{multline*}
Since
\begin{align*}
\avgint_{\partial B_{1}(0)}\left| \nabla^2\OT(x)y\right|^2\diff y=\frac{\Tr[(\nabla^2\OT(x))^2]}{
\dd}\ge \frac{(\Delta \OT(x))^2}{\dd^2},
\end{align*}
it follows that
\[
\lsh\int \left(\frac{\Delta \OT}{2\dd}\right)^p e^{-V}\diff x\ge \lc \int \frac{(\Delta \OT)^2}{\dd^2} \left(\frac{\Delta \OT}{2\dd}\right)^p e^{-V}\diff x,
\]
and hence,
\begin{align}\label{eq:p_p+2}
\int (\Delta \OT)^p e^{-V}\ge\frac{1}{\dd^2} \frac{\lc}{\lsh} \int (\Delta \OT)^{p+2} e^{-V}\diff x.
\end{align}
By H\"older's inequality, with exponents $\frac{p+2}{2},\frac{p+2}{p}$, we have
\begin{equation}\label{eq:Holder}
\begin{aligned}
\int (\Delta \OT)^p e^{-V}&\le \left(\int \left[(\Delta \OT)^p\right]^{\frac{p+2}{p}} e^{-V} \right)^{\frac{p}{p+2}}\left(\int 1^{\frac{p}{p+2}}e^{-V}\right)^{\frac{p+2}{p}}\\
&=\left(\int (\Delta\OT)^{p+2}e^{-V}\right)^{\frac{p}{p+2}},
\end{aligned}
\end{equation}
so combining~\eqref{eq:Holder} and~\eqref{eq:p_p+2} we get
\[
\left(\int (\Delta\OT)^{p+2}e^{-V}\right)^{\frac{p}{p+2}}\ge \frac{1}{\dd^2}\frac{\lc}{\lsh} \int (\Delta \OT)^{p+2} e^{-V}.
\]
Relabeling $p\mapsto 2p$, the above can be written as
\begin{align}\label{eq:q_bound_proof}
\|(\Delta\OT)^2\|_{L^{p+1}(\source)}\le\dd^2 \frac{\lsh}{\lc}.
\end{align}
Inequality~\eqref{eq:infty_bound} follows by taking $p\to\infty$,
\begin{equation}\label{eq:infty_bound_proof}
\|\Delta\OT\|_{L^{\infty}(\source)}\le \dd\sqrt{ \frac{\lsh}{\lc}}.\qedhere
\end{equation}
\end{proof}

It remains to show that the second term in~\eqref{eq:2_terms} is nonnegative, which was first observed by Kolesnikov in~\cite{MR3201654}. Here we essentially follow his proof with some minor adaptations to our case.
\begin{prop}\label{prop:noneg_term}
\begin{multline}\label{eq:noneg_term}
\int (\Delta_{\epsilon}\OT(x))^p e^{-V(x)}
\cdot\left\{\avgint_{\partial B_{\epsilon}(0)}\bigg\{ 2\langle \nabla\W(\nabla \OT(x)),\nabla \OT(x+y)-\nabla \OT(x)\rangle\right.\\
\left.\left.-\log\left[\frac{\det \nabla^2\OT(x+y)\det \nabla^2\OT(x-y)}{(\det \nabla^2\OT(x))^2}\right]\right\}\diff y\right\}\diff x\ge 0.
\end{multline}
\end{prop}
\begin{proof}
Let $\nabla\Psi=(\nabla \OT)^{-1}$ be the optimal transport map between $\diff \target=e^{-\W}\diff x$ to $\diff \source=e^{-\V}\diff x$. The key estimate required to obtain~\eqref{eq:noneg_term} is contained in the following proposition.
\begin{prop}\label{prop:nablaW_term}
\begin{multline}\label{eq:nablaW_term}
\int (\Delta_{\epsilon}\OT(x))^p e^{-V(x)}\avgint_{\partial B_{\epsilon}(0)} 2\bigl\langle \nabla\W(\nabla \OT(x)),\nabla \OT(x+y)-\nabla \OT(x)\bigr\rangle\diff y\diff x\\ 
\ge 2\avgint_{\partial B_{\epsilon}(0)}\left\{
\begin{multlined}
\int (\Delta_{\epsilon}\OT(\nabla \Psi(x)))^p\\
\left\{\Tr\left[\nabla^2\OT(\nabla\Psi(x)+y)\nabla^2\Psi(x)\right]-\dd\right\}\diff \target(x)
\end{multlined}
\right\} \diff y.
\end{multline}
\end{prop}

Assuming the estimate~\eqref{eq:nablaW_term} let us prove the estimate~\eqref{eq:noneg_term}. Using the symmetry of $\partial B_{\epsilon}(0)$ under reflection $y\mapsto -y$ we have
\begin{multline*}
\int (\Delta_{\epsilon}\OT)^p e^{-V}
\cdot\left\{\avgint_{\partial B_{\epsilon}(0)}\left\{ 2\langle \nabla\W(\nabla \OT),\nabla \OT(x+y)-\nabla \OT(x)\rangle\right.\right.\\
\shoveright{\left.\left.\ -\log\left[\frac{\det \nabla^2\OT(x+y)\det \nabla^2\OT(x-y)}{(\det \nabla^2\OT(x))^2}\right]\right\}\diff y\right\}\diff x}\\
\begin{aligned}
\overset{\eqref{eq:nablaW_term}}&{\ge} 2\avgint_{\partial B_{\epsilon}(0)}\left\{\int (\Delta_{\epsilon}\OT(\nabla \Psi))^p \left\{\Tr\left[\nabla^2\OT(\nabla\Psi(x)+y)\nabla^2\Psi(x)\right]-\dd\right\}\diff \target\right\} \diff y\\
&\quad-2\int (\Delta_{\epsilon}\OT)^p \avgint_{\partial B_{\epsilon}(0)}\log\det\left[\nabla^2\OT(x+y)(\nabla^2\OT(x))^{-1}\right]\diff y\,e^{-V}\diff x\\
&= 2\avgint_{\partial B_{\epsilon}(0)}\left\{\int \left\{\Tr\left[\nabla^2\OT(\nabla\Psi(x)+y)\nabla^2\Psi(x)\right]-\dd\right\}(\Delta_{\epsilon}\OT(\nabla \Psi))^p \diff \target\right\} \diff y\\
&\quad-2\avgint_{\partial B_{\epsilon}(0)}\left\{\int\log\det\left[\nabla^2\OT(\nabla \Psi(x)+y)(\nabla^2\OT(\nabla \Psi(x)))^{-1}\right](\Delta_{\epsilon}\OT(\nabla\Psi))^p\diff \target\right\}\diff y\\
&=2\avgint_{\partial B_{\epsilon}(0)}\left\{\left[\Tr[A(x)]-\dd-\log \det A(x)\right](\Delta_{\epsilon}\OT(\nabla\Psi))^p\diff \target\right\}\diff y,
\end{aligned}
\end{multline*}
where $A(x)=\nabla^2\OT(\nabla \Psi(x)+y)(\nabla^2\OT(\nabla \Psi(x)))^{-1}$. By~\cite[p.~8]{kolesnikov2011mass}, $\Tr[A(x)]-\dd-\log \det A(x)\ge 0$ for every $x$, which completes the proof.
\end{proof}

It remains to prove Proposition~\ref{prop:nablaW_term}.
\begin{proof}[Proof of Proposition~\ref{prop:nablaW_term}]
Let $\nabla\Psi=(\nabla \OT)^{-1}$ be the optimal transport map between $\diff \target=e^{-\W}\diff x$ to $\diff \source=e^{-\V}\diff x$. The estimate~\eqref{eq:nablaW_term} will obtained as a consequence of integration by parts. But in order to justify the vanishing of the boundary terms in the integration by parts we need to work with cutoff functions. Let $\{\eta_k\}$ be a sequence of compactly supported smooth functions $\eta_k:\R^{\dd}\to \R$ such that
\begin{itemize}
\item $0\le \eta_k\le 1$ for every $k$.
\item $\lim_{k\to\infty}\eta_k=1$ uniformly on every compact set.
\item $\lim_{k\to\infty}\int |\nabla\eta_k(x)|^2 \diff\target(x)=0$.
\end{itemize}
In the computations below we will simplify the notation and write
\[
\findiff_y(x):=\nabla \OT(x+y)-\nabla \OT(x),
\]
omitting the dependence on $x$ when it is clear from context, $\findiff_y:=\findiff_y(x)$.

By the Monge--Ampère equation~\eqref{eq:Monge_proof} (and since the integrals can be exchanged by the proof of Lemma~\ref{lem:integrate_finite}), multiplying the integrand in the left-hand side of~\eqref{eq:nablaW_term} by $\eta_k(\nabla\OT(x))$, and integrating against $e^{-\V(x)}\diff x$, yields
\begin{multline*}
\int\left\{ (\Delta_{\epsilon}\OT)^p\avgint_{\partial B_{\epsilon}(0)}2\langle \nabla\W(\nabla \OT),\findiff_y\rangle\diff y\right\}\eta_k(\nabla\OT)\diff \source\\
\begin{aligned}
&=\int\left\{ (\Delta_{\epsilon}\OT)^p\avgint_{\partial B_{\epsilon}(0)}2\langle \nabla\W(\nabla \OT),\findiff_y\rangle\diff y\right\}\eta_k(\nabla\OT)e^{-\W(\nabla\OT)}\det\nabla^2\OT\diff x\\
&=-\int\left\{ (\Delta_{\epsilon}\OT)^p\avgint_{\partial B_{\epsilon}(0)}2\left\langle \nabla\bigl(e^{-\W(\nabla\OT)}\bigr),(\nabla^2\OT)^{-1}\findiff_y\right\rangle\diff y\right\}\eta_k(\nabla\OT)\det\nabla^2\OT\diff x\\
&=-\avgint_{\partial B_{\epsilon}(0)}\left\{\int 2\left\langle \nabla\bigl(e^{-\W(\nabla\OT)}\bigr),\Cof(\nabla^2\OT)\findiff_y\right\rangle(\Delta_{\epsilon}\OT)^p\eta_k(\nabla\OT)\diff x\right\}\diff y.
\end{aligned}
\end{multline*}
By integration by parts in the $x$ variable (which has no boundary terms because $\eta_k$ is compactly supported), and using the fact that the cofactor matrix is divergence free, we get
\begin{multline*}
-\int \left\langle \nabla\bigl(e^{-\W(\nabla\OT)}\bigr),\Cof(\nabla^2\OT)\findiff_y\right\rangle(\Delta_{\epsilon}\OT)^p\eta_k(\nabla\OT)
\\
\begin{aligned}
&=\int e^{-\W(\nabla\OT)} \Div\left[\Cof(\nabla^2\OT)\findiff_y(\Delta_{\epsilon}\OT)^p\eta_k(\nabla\OT)\right]
\\
&=\int e^{-\W(\nabla\OT)} \Tr\left[\Cof(\nabla^2\OT)\nabla[\findiff_y(\Delta_{\epsilon}\OT)^p\eta_k(\nabla\OT)]\right]
\\
&=\int e^{-\W(\nabla\OT)} \Tr\left[\Cof(\nabla^2\OT)\nabla[\findiff_y(\Delta_{\epsilon}\OT)^p]\right]\eta_k(\nabla\OT)
\\
&\quad +\int e^{-\W(\nabla\OT)} \Tr\left[\Cof(\nabla^2\OT)\left\{\left((\Delta_{\epsilon}\OT)^p\findiff_y\right)\otimes\nabla[\eta_k(\nabla\OT)]\right\}\right],
\end{aligned}
\end{multline*}
where we use the convention $(\nabla u(x))_{ij}=\partial_iu_j(x)$ for a vector field $u$, and $(w\otimes v)_{ij}=w_jv_i$ for vectors $v,w$. Hence, our goal is to lower bound the term
\begin{multline}\label{eq:2terms}
\int\left\{ (\Delta_{\epsilon}\OT)^p\avgint_{\partial B_{\epsilon}(0)}2\langle \nabla\W(\nabla \OT),\findiff_y\rangle\diff y\right\}\eta_k(\nabla\OT)\diff \source\\
\begin{aligned}
&=2\avgint_{\partial B_{\epsilon}(0)}\int e^{-\W(\nabla\OT)} \Tr\left[\Cof(\nabla^2\OT)\nabla[\findiff_y(\Delta_{\epsilon}\OT)^p]\right]\eta_k(\nabla\OT)\diff x\diff y\\
&\quad+2\avgint_{\partial B_{\epsilon}(0)}\int e^{-\W(\nabla\OT)} \Tr\left[\Cof(\nabla^2\OT)\bigl\{\left((\Delta_{\epsilon}\OT)^p\findiff_y\right)\otimes\nabla[\eta_k(\nabla\OT)]\bigr\}\right]\diff x\diff y.
\end{aligned}
\end{multline}
The second term on the right-hand side of~\eqref{eq:2terms} will be shown to vanish as $k\to\infty$ so it suffices to lower bound the first term on the right-hand side of~\eqref{eq:2terms}, and then take the limit $k\to \infty$. We start with the first term in~\eqref{eq:2terms}.

\begin{lemm}\label{lem:1term}
\begin{multline}\label{eq:1term}
\avgint_{\partial B_{\epsilon}(0)}\int e^{-\W(\nabla\OT(x))}\\ \Tr\bigl[\Cof(\nabla^2\OT(x))\nabla\left[(\nabla \OT(x+y)
-\nabla \OT(x))(\Delta_{\epsilon}\OT(x))^p\right]\bigr]\eta_k(\nabla\OT(x))\diff x\diff y\\
\ge \avgint_{\partial B_{\epsilon}(0)}\int\left \{\Tr\left[\nabla^2\OT(\nabla\Psi(x)+y)\nabla^2\Psi(x)\right]-\dd\right\}\\(\Delta_{\epsilon}\OT(\nabla\Psi(x)))^p \eta_k(x)\diff \target(x)\diff y.
\end{multline}
\end{lemm}
\begin{proof}
Let $\nabla\Psi=(\nabla \OT)^{-1}$ be the optimal transport map between $\diff \target=e^{-\W}\diff x$ to $\diff \source=e^{-\V}\diff x$, and recall $\nabla^2\Psi(x)=(\nabla^2\OT)^{-1}(\nabla\Psi(x))$. Using the chain rule we have, for any vector field $u$ in $\R^{\dd}$, $\Div_x[u(\nabla \Psi(x))]=\Tr[\nabla^2 \Psi(x)\nabla u(\nabla \Psi(x))]$, so with
\begin{align*}
\findiff_y &:=\nabla \OT(x+y)-\nabla \OT(x), \\
\findiff_y(\nabla\Psi) &:=\findiff_y(\nabla\Psi(x))=\nabla \OT(\nabla\Psi(x)+y)-\nabla \OT(\nabla\Psi(x)),
\end{align*}
we have
\begin{multline*}
\int e^{-\W(\nabla\OT)} \Tr\left[\Cof\bigl(\nabla^2\OT\bigr)\nabla[(\Delta_{\epsilon}\OT)^p\findiff_y]\right]\eta_k(\nabla\OT)\diff x\\
\begin{aligned}
&=\int \Tr\left[\nabla^2\Psi(\nabla\OT)\nabla[(\Delta_{\epsilon}\OT)^p\findiff_y]\right]\eta_k(\nabla\OT)\diff \source\\
&=\int \Tr\left[\nabla^2\Psi\nabla[\findiff(\nabla\Psi)(\Delta_{\epsilon}\OT(\nabla\Psi))^p]\right]\eta_k\diff \target\\
&=\int \Div[(\Delta_{\epsilon}\OT(\nabla\Psi))^p \findiff_y(\nabla\Psi)]\eta_k\diff \target\\
&=\int \Div[\left(\nabla \OT(\nabla\Psi(x)+y)-x\right)](\Delta_{\epsilon}\OT(\nabla\Psi))^p \eta_k\diff \target\\
&\quad+\int \bigl\langle\findiff_y(\nabla\Psi),\nabla(\Delta_{\epsilon}\OT(\nabla\Psi))^p \bigr\rangle\eta_k\diff \target.
\end{aligned}
\end{multline*}
Hence,
\begin{multline}\label{eq:1term_equiv}
\avgint_{\partial B_{\epsilon}(0)}\int e^{-\W(\nabla\OT)} \Tr\left[\Cof(\nabla^2\OT)\nabla[(\Delta_{\epsilon}\OT)^p\findiff_y]\right]\eta_k(\nabla\OT)\diff x\diff y\\
=\avgint_{\partial B_{\epsilon}(0)}\int \Div[\left(\nabla \OT(\nabla\Psi(x)+y)-x\right)](\Delta_{\epsilon}\OT(\nabla\Psi))^p \eta_k\diff \target\diff y\\
+\avgint_{\partial B_{\epsilon}(0)}\int \bigl\langle\findiff_y(\nabla\Psi),\nabla(\Delta_{\epsilon}\OT(\nabla\Psi))^p \bigr\rangle\eta_k\diff \target\diff y.
\end{multline}

Let us bound the two terms in~\eqref{eq:1term_equiv} separately. For the first term in~\eqref{eq:1term_equiv} we use the chain rule to get
\begin{multline*}
\Div[\nabla \OT(\nabla\Psi(x)+y)-x]\\
=\Div[\nabla \OT(\nabla\Psi(x)+y)]-\dd\ge \Tr\left[\nabla^2\OT(\nabla\Psi(x)+y)\nabla^2\Psi(x)\right]-\dd.
\end{multline*}
Using $\nabla^2\Psi(x)=(\nabla^2\OT)^{-1}(\nabla\Psi(x))$ gives
\begin{multline}\label{eq:1term_1}
\avgint_{\partial B_{\epsilon}(0)}\int \Div[\left(\nabla \OT(\nabla\Psi(x)+y)-x\right)](\Delta_{\epsilon}\OT(\nabla\Psi))^p \eta_k\diff \target\diff y\\
\ge \avgint_{\partial B_{\epsilon}(0)}\int\left \{\Tr\left[\nabla^2\OT(\nabla\Psi(x)+y)\nabla^2\Psi(x)\right]-\dd\right\}(\Delta_{\epsilon}\OT(\nabla\Psi))^p \eta_k\diff \target\diff y.
\end{multline}
For the second term in~\eqref{eq:1term_equiv} we use the definition of $\Delta_{\epsilon}$ to write
\begin{multline}\label{eq:1term_2}
\avgint_{\partial B_{\epsilon}(0)}\int \left\langle\nabla \findiff_y(\nabla\Psi), \nabla\left((\Delta_{\epsilon}\OT(\nabla \Psi))^p\right) \right\rangle \eta_k\diff \target\diff y
\\ 
\begin{aligned}
&=p\avgint_{\partial B_{\epsilon}(0)}\int (\Delta_{\epsilon}\OT(\nabla \Psi))^{p-1} \left\langle \findiff_y(\nabla\Psi), \nabla^2\Psi \nabla\Delta_{\epsilon}\OT(\nabla \Psi)\right\rangle \eta_k\diff \target\diff y\\
&=p\int (\Delta_{\epsilon}\OT(\nabla \Psi))^{p-1} \left\langle \avgint_{\partial B_{\epsilon}(0)}\findiff_y(\nabla\Psi)\diff y,\nabla^2\Psi \nabla\Delta_{\epsilon}\OT(\nabla \Psi)\right\rangle \eta_k\diff \target \\
\overset{\eqref{eq:avg_Laplacian_def}}&{=} p\int (\Delta_{\epsilon}\OT(\nabla \Psi))^{p-1} \left\langle (\Delta_{\epsilon}\nabla \OT)(\nabla\Psi),\nabla^2\Psi \nabla\Delta_{\epsilon}\OT(\nabla \Psi)\right\rangle \eta_k\diff \target\\
&=p\int (\Delta_{\epsilon}\OT(\nabla \Psi))^{p-1} \left\langle \nabla\Delta_{\epsilon}\OT(\nabla \Psi),\nabla^2\Psi \nabla\Delta_{\epsilon}\OT(\nabla \Psi)\right\rangle \eta_k\diff \target\ge 0,
\end{aligned}
\end{multline}
where we used the convexity of $\OT$ and $\Psi$ as well as $\eta_k\ge 0$. Combining~\eqref{eq:1term_1} and~\eqref{eq:1term_2} yields~\eqref{eq:1term}.
\end{proof}

We now turn to the second term in~\eqref{eq:2terms}.
\begin{lemm}\label{lem:2term}
\begin{multline}\label{eq:2term}
2\int e^{-\W(\nabla\OT(x))} \Tr\Bigl[\Cof(\nabla^2\OT(x))\\
\bigl\{\left(\left(\nabla \OT(x+y)-\nabla \OT(x)\right)(\Delta_{\epsilon}\OT(x))^p\right)\otimes\nabla[\eta_k(\nabla\OT(x))]\bigr\}\Bigr]\diff x\\
\le 2\left(\int \bigl|\left(\nabla \OT(\nabla\Psi(x)+y)-\nabla \OT(\nabla\Psi(x))\right)(\Delta_{\epsilon}\OT(\nabla\Psi(x)))^p\bigr|^2\diff \target(x)\right)^{\frac{1}{2}}\\
\left(\int |\nabla\eta_k(x)|^2\diff \target(x)\right)^{\frac{1}{2}}.
\end{multline}
\end{lemm}
\begin{proof}
Given a symmetric matrix $M$ and vectors $u,v$ we have $\Tr[M^{-1}\{u\otimes (Mv)\}]=\langle u,v\rangle$, so applying this identity we get, with $\findiff_y:=\nabla \OT(x+y)-\nabla \OT(x)$,
\begin{multline*}
\int e^{-\W(\nabla\OT)} \Tr\left[\Cof(\nabla^2\OT)\left\{\left(\findiff_y(\Delta_{\epsilon}\OT)^p\right)\otimes\nabla[\eta_k(\nabla\OT)]\right\}\right]\diff x\\
\begin{aligned}
&=\int \Tr\bigl[\left\{\left[\findiff(\Delta_{\epsilon}\OT)^p\right]\otimes [(\nabla\eta_k)(\nabla\OT)]\right\}\bigr]\diff \source\\
&=\int \Tr\bigl[\left\{\left[\findiff_y(\nabla\Psi)(\Delta_{\epsilon}\OT(\nabla\Psi))^p\right]\right\}\otimes\nabla\eta_k\bigr]\diff \target\\
&\le \left(\int \bigl|\findiff_y(\nabla\Psi)(\Delta_{\epsilon}\OT(\nabla\Psi))^p\bigr|^2\diff \target\right)^{\frac{1}{2}}\left(\int |\nabla\eta_k|^2\diff \target\right)^{\frac{1}{2}}.\qedhere
\end{aligned}
\end{multline*}
\end{proof}

Let us now take the limits $k\to\infty$ of the two terms in~\eqref{eq:2terms}. For the first term in~\eqref{eq:2terms} we first use~\eqref{eq:1term} and then take the limit $k\to\infty$ in~\eqref{eq:1term}. We can move the limit past the integrals in~\eqref{eq:1term} by dominated convergence theorem using $\eta_k\le 1$, and the fact that the rest of the integrand is integrable (as in the proof of Lemma~\ref{lem:integrate_finite}). Thus, if the second term in~\eqref{eq:2terms} vanishes in the limit $k\to\infty$ we will get~\eqref{eq:nablaW_term}. To show that the second term in~\eqref{eq:2terms} vanishes in the limit $k\to\infty$ it suffices to show that the first term on the right-hand side of~\eqref{eq:2term} is finite since by assumption $\lim_{k\to\infty}\int |\nabla\eta_k(x)|^2 \diff\target(x)=0$. The first term on the right-hand side of~\eqref{eq:2term} is finite by the same argument as in the proof of Lemma~\ref{lem:integrate_finite}.
\end{proof}

To summarize, the combination of Lemma~\ref{lem:1term} and Lemma~\ref{lem:2term} yields Proposition~\ref{prop:nablaW_term}, which in turns implies Proposition~\ref{prop:noneg_term}. It follows that the second term in~\eqref{eq:2_terms} is nonnegative and the proof of Theorem~\ref{thm:main_OT_Lp} is complete by Lemma~\ref{lem:_term_to_thm_proof}.
\end{proof}

\begin{rema}\label{rem:modification}
Our proof of Theorem~\ref{thm:main_OT_Lp} follows the proof of~\cite[Theorem~6.2]{kolesnikov2011mass}. However, one important modification we need to make is the introduction of the $\Delta_{\epsilon}$ operator. In~\cite{kolesnikov2011mass}, the lack of regularity of $\V$ and $\nabla\OT$ is remedied by considering finite differences (rather than actual derivatives), $\V(x+e)-\V(x)$ and $\nabla\OT(x+e)-\nabla\OT(x)$, for vectors $e\in\R^{\dd}$. The finite difference $\V(x+e)-\V(x)$ is then controlled by explicit assumptions on the second directional derivatives of $\V$, which in turn leads to control on the second directional derivatives of $\OT$. In contrast, we only have at our disposal the control of $\Delta \V$, so the finite differences approach just outlined does not work. The operator $\Delta_{\epsilon}$ is the appropriate replacement to the finite differences scheme.
\end{rema}

\appendix
\section{}\label{sec:appendix}


\subsection{The operator \texorpdfstring{$\Delta_{\epsilon}$}{Delta\_ epsilon}}\label{subsec:appendix_Delta_eps}
In this section we prove some of the properties of the operator $\Delta_{\epsilon}$. In particular, let us prove Lemma~\ref{lem:avg_Laplacian}.
\begin{lemm}\label{lem:avg_Laplacian_appendix}
Fix $\epsilon> 0$ and for a function $f:\R^{\dd}\to\R$ denote
\[
\Delta_{\epsilon}f(x):=\avgint_{\partial B_{\epsilon}(0)}[f(x+y)-f(x)]\diff y,
\]
where
\[
\avgint_{\partial B_{\epsilon}(0)}f:=\frac{1}{|\partial B_{\epsilon}(0)|}\int_{\partial B_{\epsilon}(0)}f.
\]
Then,
\begin{align}\label{eq:Delta_eps_lim_appendix}
\lim_{\epsilon\to 0}\frac{\Delta_{\epsilon}f(x)}{\epsilon^{2}}=\frac{\Delta f(x)}{2\dd},
\end{align}
where $\Delta f$ is the distributional Laplacian of $f$. Further, if $\Delta f\le \ell$, then
\begin{align}\label{eq:Delta_eps_bound_appendix}
\Delta_{\epsilon}f\le \frac{\ell}{\dd}\frac{\epsilon^2}{2},\quad \forall~\epsilon>0.
\end{align}
\end{lemm}
\begin{proof}
To prove~\eqref{eq:Delta_eps_lim_appendix} first note that letting $ \nabla^2f$ be the distributional derivative of $f$ we have
\begin{align*}
\lim_{\epsilon\to 0}\frac{f(x+\epsilon y)+f(x-\epsilon y)-2f(x)}{\epsilon^2}=\left\langle \nabla^2f(x)y,y\right\rangle,
\end{align*}
in the sense of distributions, where the distribution $\langle \nabla^{2} f(x)y, y \rangle$ is defined as
\[
\left \langle \left\langle \nabla^{2} f(x)y, y \right\rangle, \eta \right \rangle =\int f(x) \left\langle \nabla^{2} \eta(x) y, y \right\rangle \diff x,
\]
for test functions $\eta$. Hence, using the symmetry of $\partial B_{\epsilon}(0)$ under reflection, and changing variables $y\mapsto \epsilon y$,
\begin{align*}
\lim_{\epsilon\to 0}\frac{\Delta_{\epsilon}f(x)}{\epsilon^2}&=\lim_{\epsilon\to 0}\frac{1}{2\epsilon^2}\left\{\avgint_{\partial B_{\epsilon}(0)}2[f(x+y)-f(x)]\diff y\right\}\\
&=\frac{1}{2}\lim_{\epsilon\to 0}\left\{\avgint_{\partial B_{\epsilon}(0)}\frac{f(x+y)+f(x-y)-2f(x)}{\epsilon^2}\diff y\right\}\\
&=\frac{1}{2}\lim_{\epsilon\to 0}\left\{\frac{1}{\epsilon^{\dd-1}}\avgint_{\partial B_{1}(0)}\frac{f(x+\epsilon y)+f(x-\epsilon y)-2f(x)}{\epsilon^2}\epsilon^{\dd-1}\diff y\right\}\\
&=\frac{1}{2}\avgint_{\partial B_{1}(0)}\left\langle \nabla^2f(x)y,y\right\rangle \diff y.
\end{align*}
For each $i,j\in [\dd]$, by the symmetry of $\partial B_{\epsilon}(0)$ under $y_j\mapsto -y_j$,
\[
\avgint_{\partial B_{1}(0)}y_iy_j \diff y=\delta_{ij}\avgint_{\partial B_{1}(0)}y_i^2 \diff y=\frac{1}{\dd}.
\]
It follows that
\begin{align*}
\lim_{\epsilon\to 0}\frac{\Delta_{\epsilon}f(x)}{\epsilon^2}&=\frac{1}{2}\sum_{i,j=1}^{\dd}\avgint_{\partial B_{1}(0)}\partial_{ij}^2f(x)y_iy_j \diff y=\frac{1}{2}\sum_{i=1}^{\dd}\partial_{ii}^2f(x)\avgint_{\partial B_{1}(0)}y_i^2\diff y\\
&=\frac{1}{2\dd}\sum_{i=1}^{\dd}\partial_{ii}^2f(x)=\frac{\Delta f(x)}{2\dd}.
\end{align*}



Next we move to the proof of~\eqref{eq:Delta_eps_bound_appendix}. We have
\begin{align*}
\Delta_{\epsilon}f(x)&=\int_0^{\epsilon} \frac{\diff}{\diff r}\left\{\avgint_{\partial B_{r}(0)}[f(x+y)-f(x)]\diff y\right\}\diff r\\
&=\int_0^{\epsilon} \frac{\diff}{\diff r}\left[\avgint_{\partial B_{r}(0)}f(x+y)\diff y\right]\diff r.
\end{align*}
Since, by the change of variables $y\mapsto r y$,
\begin{align*}
\avgint_{\partial B_{r}(0)}f(x+y)\diff y=\avgint_{\partial B_1(0)}f(x+ry)\diff y,
\end{align*}
we have, by integration by parts,
\begin{align*}
&\frac{\diff}{\diff r}\left[\avgint_{\partial B_{r}(0)}f(x+y)\diff y\right]=\avgint_{\partial B_1(0)}\frac{\diff}{\diff r}f(x+ry)\diff y=\avgint_{\partial B_1(0)}\langle\nabla f(x+ry),y\rangle\diff y\\
&=\avgint_{\partial B_1(0)}\frac{1}{r}\langle\nabla_y [f(x+ry)],y\rangle\diff y=\avgint_{ B_1(0)}\frac{1}{r}\Delta_y[f(x+ry)]\diff y\\
&=\frac{r}{\dd}\avgint_{B_1(0)}\Delta f(x+ry)\diff y\le \frac{\ell r}{\dd},
\end{align*}
where the last inequality used $\Delta f\le \ell$. It follows that
\[
\Delta_{\epsilon}f(x)\le \int_0^{\epsilon}r\diff r=\frac{\ell}{\dd}\frac{\epsilon^2}{2}.\qedhere
\]
\end{proof}

\subsection{Smoothing under Ornstein--Uhlenbeck and heat semigroups}\label{subsec:appendix_smooth}
In this section we discuss Proposition~\ref{prop:OU_smooth}. Since most of the results in the proposition are standard, when relevant we will only sketch the arguments and give references for detailed proofs. For the sake of completeness we will prove some additional smoothing properties beyond those stated in Proposition~\ref{prop:OU_smooth}.



\begin{prop}\label{prop:OU_smooth_appendix}
Let $\Gaussian$ be the standard Gaussian measure on $\R^{\dd}$, and let $f:\R^{\dd}\to \R_{\ge 0}$ be a nonnegative function in $L^1(\Gaussian)$. Let $(\Pheat_t)_{t\ge 0}$ be the Ornstein--Uhlenbeck semigroup,
\begin{equation}\label{eq:OU_def_prop_appendix}
\Pheat_t\density(x):=\int_{\R^{\dd}} \density\left(e^{-t}x+\sqrt{1-e^{-2t}}y\right)\diff \Gaussian(y),\quad t\ge 0,\quad x\in \R^{\dd},
\end{equation}
and let $(\Hheat_t)_{t\ge 0}$ be the heat semigroup
\begin{equation}\label{eq:heat_def_prop_appendix}
\Hheat_t\density(x):=\int_{\R^{\dd}} \density\left(x+\sqrt{t}y\right)\diff \Gaussian(y),\quad t\ge 0,\quad x\in \R^{\dd}.
\end{equation}
\begin{enumerate}
\item\label{propA.2.1} For any $x\in \R^{\dd}$ and $t>0$,
\begin{equation}\label{eq:less_lc_appendix}
\nabla^2\log\Pheat_t\density(x) \succeq -\frac{e^{-2t}}{1-e^{-2t}} \Id_{\dd}\quad\text{and}\quad \nabla^2\log\Hheat_t\density(x) \succeq -\frac{1}{t}\Id_{\dd}.
\end{equation}
\item\label{propA.2.2} Suppose that $\density$ is $\conv$-log-concave for some $\conv\in \R$ (i.e., $x\mapsto \log \density(x)-\conv\frac{|x|^2}{2}$ is concave). Then, for every $x\in \R^{\dd}$,
\begin{equation}\label{eq:beta_log-concave_OU_appendix}
\nabla^2\log\Pheat_t\density(x)\preceq \frac{ e^{-2t}\conv}{1-\conv(1-e^{-2t})} \Id_{\dd}
\begin{cases}
\text{for any }t\in [0,\infty) &\text{if }\conv\le 1\\
\text{for any }t\in \left[0,\log\left(\sqrt{\frac{\conv}{\conv-1}}\right)\right] &\text{if }\conv >1
\end{cases},
\end{equation}
and
\begin{equation}\label{eq:beta_log-concave_heat_appendix}
\nabla^2\log\Hheat_t\density(x)\preceq \frac{\conv}{1-\conv t}\Id_{\dd}
\begin{cases}
\text{for any }t\in [0,\infty) &\text{if }\conv\le 0\\
\text{for any }t\in \left[0,\frac{1}{\conv}\right] &\text{if }\conv >0
\end{cases}.
\end{equation}

\item\label{propA.2.3} Suppose that $\density$ is $\conv$-log-convex for some $\conv\le 0$ (i.e., $x\mapsto \log \density(x)-\conv\frac{|x|^2}{2}$ is convex). Then, for every $x\in \R^{\dd}$ and $t>0$,
\begin{equation}\label{eq:beta_cvx_appendix}
\nabla^2\log\Pheat_t\density(x)\succeq \frac{e^{-2t}\conv}{1-\conv(1-e^{-2t})} \Id_{\dd}\quad\text{and}\quad \nabla^2\log\Hheat_t\density(x) \succeq \frac{\conv}{1-\conv t}\Id_{\dd}.
\end{equation}
Suppose $\density$ is $\conv\dd$-log-subharmonic for some $\conv\le 0$ (i.e., $x\mapsto \log f(x)-\conv\frac{|x|^2}{2}$ is subharmonic). Then, for every $x\in \R^{\dd}$ and $t>0$,
\begin{equation}\label{eq:beta_subharmonic_appendix}
\Delta\log\Pheat_t\density(x)\ge \frac{e^{-2t}\conv\dd}{1-\conv(1-e^{-2t})}\quad\text{and}\quad \Delta \log\Hheat_t\density(x) \ge \frac{\conv\dd}{1-\conv t}.
\end{equation}
\end{enumerate}
\end{prop}
\begin{rema}
It was shown in~\cite[Remark]{MR4632741} that there exists a $0$-log-superharmonic $\density$ such that, for any $t>0$, $\Pheat_t\density$ is strictly log-subharmonic. Concretely, let $\dd=2$ and $\density(x)=\density(x_1,x_2):=e^{x_1x_2}$. Then, $\Delta\log
\density(x)=0$ for all $x\in \R^{\dd}$ while $\Delta\log
\Pheat_t\density(x)>0$ for all $x\in\R^{\dd}$ and $t>0$. We conclude that log-superharmonic functions are not preserved under the Ornstein--Uhlenbeck/heat semigroups.
\end{rema}



\begin{proof}
We will present a number of proofs for the various parts in Proposition~\ref{prop:OU_smooth_appendix}. Since the Ornstein--Uhlenbeck semigroup and heat semigroup are related by the change of variables
\begin{equation}\label{eq:OU_heat}
\Pheat_t\density(x)=\Hheat_{1-e^{-2t}}\density (e^{-t}x),
\end{equation}
we will use whichever one is convenient for the specific proof. The most robust proof technique (which works also on manifolds) is due to Hamilton~\cite{MR4370325,hamilton1993matrix} based on deriving a partial differential equation for $M(t,x):=\nabla^2\log\Hheat_t\density(x)$, and then using the maximum principle. Standard computation (e.g.~\cite[Proposition~2.11]{muller2006differential}) shows that
\begin{equation}\label{eq:PDE_Hessian}
\partial_tM(t,x)=\frac{1}{2}\Delta M(t,x)+L(M(t,x))+M^2(t,x),
\end{equation}
where $L(M(t,x))$ is a \emph{linear} differential operator in $M(t,x)$. Below we will apply the maximum principle to $(t,x)\mapsto M(t,x)$ to deduce Proposition~\ref{prop:OU_smooth_appendix}$\MK$(1-3).

Another proof technique, which is more probabilistic in nature, uses the specific from of the Ornstein--Uhlenbeck/heat semigroup in Euclidean space. Specifically, it is based on the following covariance identity
\begin{equation}\label{eq:cov_identity}
\nabla^2\log\Pheat_t\density(x)=\frac{e^{-2t}}{1-e^{-2t}}\left(\frac{\Cov\left[p_{e^{-t}x,1-e^{-2t}}\right]}{1-e^{-2t}}-\Id_{\dd}\right), \quad \forall~x\in\R^{\dd},~\forall~t>0,
\end{equation}
where $\Cov\left[p_{z,s}\right]$ is the covariance matrix of the probability measure
\begin{equation}\label{eq:conditional_measure}
\diff p_{z,s}(y)\propto \density(y)e^{-\frac{|y-z|^2}{2s}}\diff y.
\end{equation}
The identity~\eqref{eq:cov_identity} is standard, e.g., \cite[Equation (3.2)]{MR4797372}, \cite[Equation (3.4)]{MR4748461}.


\begin{proofc}[{Proofs of Proposition~\ref{prop:OU_smooth_appendix}$\MK$\eqref{propA.2.1}}]
\begin{proof}[{Proof 1}]
This is Hamilton's matrix inequality~\cite{MR4370325} which is proven via the maximum principle. Indeed, by~\eqref{eq:PDE_Hessian}, since the smallest eigenvalue \(\lambda\) is a concave function of \(M(x,t)\), it satisfies
\begin{equation}\label{eq:min_eigenvalue}
\partial_{t}\lambda(x,t)\ge \frac{1}{2}\Delta \lambda (x,t) +L(\lambda (x,t))+ \lambda^{2}(x,t).
\end{equation}
Hence, $\lambda(x,t)\ge g(t)$ where \(g(t):=-1/t\) solves the equation,
\[
\partial_tg(t)=g^{2}(t), \qquad g(0)=-\infty.
\]
\let\qed\relax
\end{proof}
\begin{proof}[{Proof 2}] Use~\eqref{eq:cov_identity} and $\Cov\left[p_{e^{-t}x,e^{-2t}}\right]\succeq 0$.
\let\qed\relax
\end{proof}
\begin{proof}[{Proof 3}]
Consequence of the intrinsic dimensional local logarithmic Sobolev inequality: the term inside the logarithm in~\cite[Eq.~29]{MR4698563} must be nonnegative. (This is analogous to the way in which the Li--Yau inequality is deduced from the dimensional local logarithmic Sobolev inequality~\cite{MR3612336}.)
\let\qed\relax
\end{proof}

\begin{proof}[{Proof 4}]
Mixture of log-convex densities is log-convex: for any fixed $y\in \R^{\dd}$ the function $x\mapsto \frac{e^{-2t}}{1-e^{-2t}}\frac{|x|^2}{2}+\log \Pheat_t\delta_y(x)$ is convex, where $\delta_y$ is a point mass at $y$. Hence, the mixture $x\mapsto\frac{e^{-2t}}{1-e^{-2t}}\frac{|x|^2}{2}+\log \int\density(y)\Pheat_t\delta_y(x)\diff y$ is also convex for any probability density $\density:\R^{\dd}\to \R_{\ge 0}$. The proof is complete since $ \int\density(y)\Pheat_t\delta_y(x)\diff y=\Pheat_t\density(x)$~\cite[Lemma~1.3, Appendix]{MR3782065}.\qedhere
\end{proof}\let\qed\relax
\end{proofc}

\begin{proofc}[{Proofs of Proposition~\ref{prop:OU_smooth_appendix}$\MK$\eqref{propA.2.2}}]
\begin{proof}[{Proof 1}]
We will use the maximum principle. Let $\Lambda(t,x)$ be the maximal eigenvalue of $M(t,x):=\nabla^2\log\Hheat_t\density(x)$. Since the maximal eigenvalue $\Lambda$ is a convex function of $M$, \eqref{eq:PDE_Hessian} implies that
\begin{equation}
\partial_t\Lambda(t,x)\le \frac{1}{2}\Delta \Lambda(t,x)+L(\Lambda(t,x))+\Lambda^2(t,x).
\end{equation}
Fix $x\in \R^{\dd}$. Then,
\begin{equation*}
\Lambda(t,x)\le g(t)\quad\text{where $g$ solves the equation} \quad \partial_tg(t)=g^2(t),\quad g(0)=\Lambda(0,x),
\end{equation*}
for every $t$ for which $g$ is well-defined. Since the solution of the ordinary differential equation for $g$ is $g(t)=\frac{\Lambda(0,x)}{1-\Lambda(0,x)t}$ for all $t$ where the denominator does not vanish, we~see that if $\Lambda(0,x)\le 0$, then $g$ is well-defined for all $t\ge 0$, and if $\Lambda(0,x)> 0$, then $g$ is well-defined for all $t\in \left[0,\frac{1}{\Lambda(0,x)}\right)$. Hence, since $\Lambda(0,x)\le \conv$,
\begin{equation*}
\Lambda(t,x)\le\frac{\Lambda(0,x)}{1-\Lambda(0,x)t}
\begin{cases}
\text{for any }t\in [0,\infty) &\text{if }\conv\le 0\\
\text{for any }t\in \left[0,\frac{1}{\conv}\right] &\text{if }\conv >0,
\end{cases}
\end{equation*}
which completes the proof of~\eqref{eq:beta_log-concave_heat_appendix}.\let\qed\relax
\end{proof}

\begin{proof}[{Proof 2}]
We will use the covariance identity following the argument in~\cite[Lemma~3.4$\MK$(2)]{MR4797372}. If $\density$ is $\conv$-log-concave then the measure $p_{z,s}$ in~\eqref{eq:conditional_measure} is $(\frac{1}{s}-\conv)$-log-concave. In particular, $p_{e^{-t}x,1-e^{-2t}}$ is $(\frac{1}{1-e^{-2t}}-\conv)$-log-concave so by the Brascamp--Lieb inequality~\cite{MR450480}, as long as $(\frac{1}{1-e^{-2t}}-\conv)\ge 0$,
\[
\Cov\left[p_{e^{-t}x,1-e^{-2t}}\right] \preceq \left(\frac{1}{1-e^{-2t}}-\conv\right)^{-1}\Id_{\dd}.
\]
Hence, by~\eqref{eq:cov_identity},
\begin{equation*}
\nabla^2\log\Pheat_t\density(x)\preceq \frac{\conv e^{-2t}}{1-\conv(1-e^{-2t})} \Id_{\dd}
\begin{cases}
\text{for any }t\in [0,\infty) &\text{if }\conv\le 1\\
\text{for any }t\in \left[0,\log\left(\sqrt{\frac{\conv}{\conv-1}}\right)\right] &\text{if }\conv >1,
\end{cases}
\end{equation*}
which proves~\eqref{eq:beta_log-concave_OU_appendix}.\let\qed\relax
\end{proof}

\begin{proof}[{Proof 3}]
We will use the Prékopa--Leindler inequality. We need to show that $x\mapsto\Hheat_t\density(x)e^{-c_t\frac{|x|^2}{2}}$ is log-concave were $c_t:=\frac{\conv}{1-\conv t}$. We have
\begin{align*}
&\Hheat_t\density(x)e^{-\frac{\conv}{1-\conv t}\frac{|x|^2}{2}}=\int \density(z)e^{-\frac{|z-x|^2}{2t}}e^{-\frac{\conv}{1-\conv t}\frac{|x|^2}{2}}\diff z\\
&=\int \density(z)\exp\left(-\frac{1}{2}\left[\frac{1}{t(1+tc_t)}\left|z-(1+tc_t)x\right|^2+\frac{c_t}{1+tc_t}|z|^2\right]\right)\diff z\\
&=\Hheat_{t(1+tc_t)}\density_{c_t}((1+tc_t)x),
\end{align*}
where $\density_{c_t}(z)=\density(z)e^{-\frac{1}{2}\frac{c_t}{1+tc_t}|z|^2}$, as long as $1+tc_t\ge 0$. Since $\density$ is $\conv$-log-concave, the function $\density_{c_t}$ is $0$-log-concave as $\conv-\frac{c_t}{1+tc_t}=\conv-\frac{\frac{\conv}{1-\conv t}}{1+t\frac{\conv}{1-\conv t}}=0$. Hence, $\Hheat_{t(1+tc_t)}\density_{c_t}$ is log-concave because it is the convolution of log-concave functions, which is log-concave by the Prékopa--Leindler inequality~\cite[\S~9]{gardner2002brunn}. Hence, as long as $1+tc_t\ge 0$, $\Hheat_t\density(x)e^{-\frac{\conv}{1-\conv t}\frac{|x|^2}{2}}$ is log-concave, which implies~\eqref{eq:beta_log-concave_heat_appendix}.
\end{proof}\let\qed\relax
\end{proofc}

\begin{proof}[{Proofs of Proposition~\ref{prop:OU_smooth_appendix}$\MK$\eqref{propA.2.3}}]

\begin{proof}[{Proof 1}]
We will use the maximum principle. Recall that the minimum eigenvalue $\lambda(t,x)$ of $M(t,x):=\nabla^2\log\Hheat_t\density(x)$ satisfies~\eqref{eq:min_eigenvalue}. Hence, $\lambda(t,x)\ge g(t)$ where $g$ solves the equation $\partial_tg(t)=g^2(t)$ with $g(0)=\lambda(0,x)$. Since the solution of the ordinary differential equation for $g$ is $g(t)=\frac{\lambda(0,x)}{1-\lambda(0,x)t}$ for all $t$ where the denominator does not vanish, we get that if $\lambda(0,x)\le 0$, then, for all $t\ge 0$, $\lambda(t,x)\ge\frac{\lambda(0,x)}{1-\lambda(0,x)t}$. Since $\lambda(0,x)\ge \conv$ the proof of~\eqref{eq:beta_cvx_appendix} is complete.


To prove~\eqref{eq:beta_subharmonic_appendix} let $m(t,x):=\Tr[M(t,x)]$, and note that, by $\Tr[M^2(t,x)]\ge \frac{\Tr[M(t,x)]^2}{\dd}$, \eqref{eq:PDE_Hessian} implies that
\begin{equation}
\partial_tm(t,x)\ge \frac{1}{2}\Delta m(t,x)+L(m(t,x))+\frac{m^2(t,x)}{\dd}.
\end{equation}
Again, this implies that, for any fixed $x\in \R^{\dd}$ and $t\ge 0$, $m(t,x)\ge \frac{m(0,x)}{1-\frac{m(0,x)}{\dd}t}$. Since $m(t,x)\ge\conv\dd$ the proof of~\eqref{eq:beta_subharmonic_appendix} is complete.\let\qed\relax
\end{proof}

\begin{proof}[{Proof 2}]
The proof is based on the fact that a mixture of log-convex (res. log-subharmonic) functions is log-convex (res. log-subharmonic): the argument we present below is based on the proof of~\cite[Lemma~5]{MS2022} which treats the $\conv$-log-convex case. Here we adapt the proof also to the $\conv\dd$-log-subharmonic case. Our starting point is the following identity for the action of Ornstein--Uhlenbeck semigroup. The point of this identity to separate the effect of the semigroup on the quadratic part of the function.
\begin{lemm}\label{lem:Ptf_identity}
Fix $\beta\ge 0$ and let $\density(x)=e^{-R(x)-\beta\frac{|x|^2}{2}}$. Let $(\Pheat_t)$ be the Ornstein--Uhlenbeck semigroup~\eqref{eq:OU_def_prop_appendix}. Then,
\begin{align}\label{eq:P_t_rep_quadratic}
\Pheat_t\density(x)=e^{-\frac{\beta e^{-2t}}{1+(1-e^{-2t})\beta}\frac{|x|^2}{2}}\int e^{-R_{t,z}(x)}e^{-\frac{\beta}{2}|z|^2}\diff z,
\end{align}
where
\begin{multline}\label{eq:R_tz_def}
R_{t,z}(x):=(2\pi)^{-\frac{\dd}{2}}\left(1-e^{-2t}+\beta^{-1}\right)^{-\frac{\dd}{2}}\\
R\left(\frac{\sqrt{1-e^{-2t}}}{\sqrt{1-e^{-2t}+\beta^{-1}}}z+\frac{e^{-t}}{1+\beta(1-e^{-2t})}x\right).
\end{multline}
\end{lemm}
\begin{proof}
By definition,
\begin{align*}
&\Pheat_t\density(x)\\
&=\int \density\left(e^{-t}x+\sqrt{1-e^{-2t}}y\right)\diff \Gaussian(y)\\
&=\frac{1}{(2\pi)^{\frac{\dd}{2}}}\int e^{-R\left(e^{-t}x+\sqrt{1-e^{-2t}}y\right)-\beta\frac{\left|e^{-t}x+\sqrt{1-e^{-2t}}y\right|^2}{2}-\frac{|y|^2}{2}}\diff y\\
&=\frac{e^{-\beta e^{-2t}\frac{|x|^2}{2}}}{(2\pi)^{\frac{\dd}{2}}}\int e^{-R\left(e^{-t}x+\sqrt{1-e^{-2t}}y\right)-\beta\frac{2\left\langle e^{-t}x,\sqrt{1-e^{-2t}}y\right\rangle +\left(1-e^{-2t}+\beta^{-1}\right)|y|^2}{2}}\diff y\\
&=\frac{e^{-\beta e^{-2t}\frac{|x|^2}{2}}}{(2\pi)^{\frac{\dd}{2}}}\int e^{-R\left(e^{-t}x+\sqrt{1-e^{-2t}}y\right)-\frac{\beta}{2}\left|\sqrt{1-e^{-2t}+\beta^{-1}}y+\frac{\sqrt{1-e^{-2t}}e^{-t}}{\sqrt{1-e^{-2t}+\beta^{-1}}}x\right|^2}\!\!e^{-\frac{-\left(1-e^{-2t}\right)e^{-2t}}{1-e^{-2t}+\beta^{-1}}\frac{\beta|x|^2}{2}}\diff y\\
&=\frac{e^{-\left(1-\frac{1-e^{-2t}}{1-e^{-2t}+\beta^{-1}}\right)\frac{e^{-2t}\beta}{2}|x|^2}}{(2\pi)^{\frac{\dd}{2}}}\int e^{-R\left(e^{-t}x+\sqrt{1-e^{-2t}}y\right)-\frac{\beta}{2}\left|\sqrt{1-e^{-2t}+\beta^{-1}}y+\frac{\sqrt{1-e^{-2t}}e^{-t}}{\sqrt{1-e^{-2t}+\beta^{-1}}}x\right|^2}\diff y\\
&=\frac{e^{-\frac{\beta e^{-2t}}{1+\left(1-e^{-2t}\right)\beta}\frac{|x|^2}{2}}}{(2\pi)^{\frac{\dd}{2}}}\int e^{-R\left(e^{-t}x+\sqrt{1-e^{-2t}}y\right)-\frac{\beta}{2}\left|\sqrt{1-e^{-2t}+\conv^{-1}}y+\frac{\sqrt{1-e^{-2t}}e^{-t}}{\sqrt{1-e^{-2t}+\beta^{-1}}}x\right|^2}\diff y.
\end{align*}

Let $z=\sqrt{1-e^{-2t}+\beta^{-1}}y+\frac{\sqrt{1-e^{-2t}}e^{-t}}{\sqrt{1-e^{-2t}+\beta^{-1}}}x$, so that $\diff y=(1-e^{-2t}+\beta^{-1})^{-\frac{\dd}{2}}\diff z$, and change variables to get
\begin{multline*}
\Pheat_t\density(x)\\
\begin{aligned}
&=\frac{e^{-\frac{\beta e^{-2t}}{1+\left(1-e^{-2t}\right)\beta}\frac{|x|^2}{2}}}{(2\pi)^{\frac{\dd}{2}}}\int e^{-R\left(\frac{\sqrt{1-e^{-2t}}}{\sqrt{1-e^{-2t}+\beta^{-1}}}z+\frac{e^{-t}}{1+\beta\left(1-e^{-2t}\right)}x\right)-\frac{\beta}{2}|z|^2}\left(1-e^{-2t}+\beta^{-1}\right)^{-\frac{\dd}{2}}\diff z\\
&=e^{-\frac{\beta e^{-2t}}{1+\left(1-e^{-2t}\right)\beta}\frac{|x|^2}{2}}\int e^{-R_{t,z}(x)}e^{-\frac{\beta}{2}|z|^2}\diff z,
\end{aligned}
\end{multline*} 
where
\begin{multline*}
R_{t,z}(x)\\
:=(2\pi)^{-\frac{\dd}{2}}\left(1-e^{-2t}+\beta^{-1}\right)^{-\frac{\dd}{2}}R\left(\frac{\sqrt{1-e^{-2t}}}{\sqrt{1-e^{-2t}+\beta^{-1}}}z+\frac{e^{-t}}{1+\beta\left(1-e^{-2t}\right)}x\right).\qedhere
\end{multline*}
\end{proof}

Let us now complete the proof of Proposition~\ref{prop:OU_smooth_appendix}$\MK$\eqref{propA.2.3} using~\eqref{eq:P_t_rep_quadratic}. Given $\density$ which is $\conv$-log-convex (res. $\conv\dd$-log-subharmonic) let $R(x):=-\log\density (x)-|\conv|\frac{|x|^2}{2}$. Then $R$ is concave (res. superharmonic), so with $R_{t,z}$ as in~\eqref{eq:R_tz_def} we have that $x\mapsto e^{-R_{t,z}(x)}$ is log-convex (res. log-subharmonic). Since the mixture of log-convex functions is log convex~\cite[Chapter~16.B]{marshall2011inequalities} (res. the mixture of log-subharmonic functions is log-subharmonic~\cite[Proposition~2.2]{MR2578455}), it follows that $x\mapsto \int e^{-R_{t,z}(x)}e^{-\frac{|\conv|}{2}|z|^2}\diff z$ is log-convex (res. log-subharmonic). From~\eqref{eq:P_t_rep_quadratic} we see that $\Pheat_t\density$ is $\frac{\conv e^{-2t}}{1-\conv(1-e^{-2t})}$-log-convex (res. $\dd\frac{\conv e^{-2t}}{1-\conv(1-e^{-2t})}$-log-subharmonic).
\end{proof}\let\qed\relax
\end{proof}\let\qed\relax
\end{proof}

\begin{thebibliography}{99}
\bibitem[Caf00]{MR1800860} L.A.Caffarelli, ``Monotonicity properties of optimal transportation and the FKG and related inequalities'', \emph{Commun. Math. Phys.}\textbf{214}(2000),no.3,547--563.
\bibitem[CEF19]{Dario2019} D.Cordero-Erausquin and A.Figalli, ``Regularity of monotone transport maps between unbounded domains'', \emph{Discrete Contin. Dyn. Syst.}\textbf{39}(2019),no.12,7101--7112.
\bibitem[Kol11]{kolesnikov2011mass} A.V.Kolesnikov, ``Mass transportation and contractions'',2011, arXiv:1103.1479.
\bibitem[Kol13]{MR3201654} A.V.Kolesnikov, ``On Sobolev regularity of mass transport and transportation inequalities'', \emph{Theory Probab. Appl.}\textbf{57}(2013),no.2,243--264.
\bibitem[MS23]{MS2022} D.Mikulincer and Y.Shenfeld, ``On the Lipschitz Properties of Transportation Along Heat Flows'', \emph{Geometric Aspects of Functional Analysis}, Lect. Notes in Math.2327, Springer,2023,269--290.
\bibitem[Nee22]{neeman2022lipschitz} J.Neeman, ``Lipschitz changes of variables via heat flow'',2022, arXiv:2201.03403.
\bibitem[KM12]{MR2983070} Y.-H.Kim and E.Milman, ``A generalization of Caffarelli's contraction theorem via (reverse) heat flow'', \emph{Math. Ann.}\textbf{354}(2012),no.3,827--862.
\bibitem[ACCM23]{MR4370325} G.Ascione,D.Castorina,G.Catino and C.Mantegazza, ``A matrix Harnack inequality for semilinear heat equations'', \emph{Math. Eng.}\textbf{5}(2023),no.1,article003(15pages).
\bibitem[Ham93]{hamilton1993matrix} R.S.Hamilton, ``Matrix Harnack estimate for the heat equation'', \emph{Commun. Anal. Geom.}\textbf{1}(1993),no.1,113--126.
\bibitem[Mül06]{muller2006differential} R.Müller, \emph{Differential Harnack inequalities and the Ricci flow}, EMS Series of Lectures in Mathematics5, European Mathematical Society,2006.
\bibitem[KP23]{MR4748461} B.Klartag and E.Putterman, ``Spectral monotonicity under Gaussian convolution'', \emph{Ann. Fac. Sci. Toulouse Math.}(6)\textbf{32}(2023),no.5,939--967.
\bibitem[MS24]{MR4797372} D.Mikulincer and Y.Shenfeld, ``The Brownian transport map'', \emph{Probab. Theory Relat. Fields}\textbf{190}(2024),no.1--2,379--444.
\bibitem[ES24]{MR4698563} A.Eskenazis and Y.Shenfeld, ``Intrinsic dimensional functional inequalities on model spaces'', \emph{J. Funct. Anal.}\textbf{286}(2024),no.7,article110338(56pages).
\bibitem[BBG17]{MR3612336} D.Bakry,F.Bolley and I.Gentil, ``The Li--Yau inequality and applications under a curvature-dimension condition'', \emph{Ann. Inst. Fourier}\textbf{67}(2017),no.1,397--421.
\bibitem[EL18]{MR3782065} R.Eldan and J.R.Lee, ``Regularization under diffusion and anticoncentration of the information content'', \emph{Duke Math. J.}\textbf{167}(2018),no.5,969--993.
\bibitem[BL76]{MR450480} H.J.Brascamp and E.H.Lieb, ``On extensions of the Brunn--Minkowski and Prékopa--Leindler theorems,including inequalities for log concave functions,and with an application to the diffusion equation'', \emph{J. Funct. Anal.}\textbf{22}(1976),no.4,366--389.
\bibitem[Gar02]{gardner2002brunn} R.Gardner, ``The Brunn--Minkowski inequality'', \emph{Bull. Am. Math. Soc.}\textbf{39}(2002),no.3,355--405.
\bibitem[BNT23]{MR4632741} N.Bez,S.Nakamura and H.Tsuji, ``Stability of hypercontractivity,the logarithmic Sobolev inequality,and Talagrand's cost inequality'', \emph{J. Funct. Anal.}\textbf{285}(2023),no.10,article110121(66pages).
\bibitem[MOA11]{marshall2011inequalities} A.W.Marshall,I.Olkin and B.C.Arnold, \emph{Inequalities: theory of majorization and its applications}, second ed., Springer Series in Statistics, Springer,2011.
\bibitem[GKL10]{MR2578455} P.Graczyk,T.Kemp and J.-J.Loeb, ``Hypercontractivity for log-subharmonic functions'', \emph{J. Funct. Anal.}\textbf{258}(2010),no.6,1785--1805.
\end{thebibliography}
\end{document}
